\documentclass[conference]{IEEEtran}

\usepackage{amsmath}
\usepackage{amssymb}
\usepackage{amsthm}
\usepackage{graphicx}
\graphicspath{{figures/standalone/}{figures/}{schematics/}{.}}
\usepackage{booktabs}
\usepackage{multirow}
\usepackage{array}
\usepackage{algorithm}
\usepackage{algorithmic}
\usepackage{cite}
\usepackage{url}

% Every empirical number in the text is \LV{key}, generated from the live benchmark CSVs by
% benchmarks/scripts/make_paper_tables_live.py; an unknown key stops compilation.
\makeatletter
\newcommand{\LVdef}[2]{\expandafter\def\csname LV@#1\endcsname{#2}}
\newcommand{\LV}[1]{\ifcsname LV@#1\endcsname\csname LV@#1\endcsname\else\errmessage{Undefined live value #1}\fi}
\makeatother
\input{live/live_values.tex}

\setlength{\textfloatsep}{3pt plus 1pt minus 1pt}
\setlength{\floatsep}{3pt plus 1pt minus 1pt}
\setlength{\intextsep}{3pt plus 1pt minus 1pt}
\setlength{\dbltextfloatsep}{3pt plus 1pt minus 1pt}
\setlength{\dblfloatsep}{3pt plus 1pt minus 1pt}
\setlength{\abovedisplayskip}{3pt plus 1pt minus 1pt}
\setlength{\belowdisplayskip}{3pt plus 1pt minus 1pt}
\setlength{\abovedisplayshortskip}{1pt plus 1pt}
\setlength{\belowdisplayshortskip}{2pt plus 1pt minus 1pt}
\renewcommand{\topfraction}{0.95}
\renewcommand{\bottomfraction}{0.95}
\renewcommand{\textfraction}{0.05}
\renewcommand{\floatpagefraction}{0.85}
\renewcommand{\dbltopfraction}{0.95}
\renewcommand{\dblfloatpagefraction}{0.85}

% Theorem Environments
\newtheorem{proposition}{Proposition}
\newtheorem{theorem}{Theorem}
\newtheorem{lemma}{Lemma}
\newtheorem{corollary}{Corollary}
\newtheorem{remark}{Remark}
\newtheorem{definition}{Definition}

\makeatletter
\def\thebibliography#1{\section*{\refname}%
    \fontsize{5.7pt}{6.7pt}\selectfont
    \list{\@biblabel{\@arabic\c@enumiv}}%
    {\settowidth\labelwidth{\@biblabel{#1}}%
    \leftmargin\labelwidth
    \advance\leftmargin\labelsep
    \setlength{\itemsep}{-0.4pt}%
    \setlength{\parsep}{0pt}%
    \setlength{\topsep}{0pt}%
    \usecounter{enumiv}%
    \let\p@enumiv\@empty
    \renewcommand\theenumiv{\@arabic\c@enumiv}}%
    \let\newblock\@empty
    \sloppy\clubpenalty4000\widowpenalty4000%
    \sfcode`\.=1000\relax}
\makeatother

\begin{document}

% =========================================================
% TITLE
% =========================================================

\title{
Fast-WCS: A Hierarchical 2D-to-3D Analytical Optimization Framework for Attributed Graph Matching and Closed-Form Geometric Superposition
}

% =========================================================
% AUTHOR BLOCK -- 3 + 2 SYMMETRIC LAYOUT (EXACTLY 5 AUTHORS)
% =========================================================

\author{
\begingroup
\renewcommand{\arraystretch}{0.90}
\setlength{\tabcolsep}{2pt}

\begin{tabular}{@{}ccc@{}}

% ROW 1: AUTHORS 1, 2, 3
\begin{minipage}[t]{0.30\textwidth}
\centering
\fontsize{8}{9}\selectfont
\textsuperscript{1st} Dr. T Ramraj

\vspace{1pt}

\textit{Department of Computer Science}\\
\textit{Amrita Vishwa Vidyapeetham}\\
Coimbatore, India\\
t\_ramraj@cb.amrita.edu
\end{minipage}
&
\begin{minipage}[t]{0.34\textwidth}
\centering
\fontsize{8}{9}\selectfont
\textsuperscript{2nd} Allampati Chenchu Sanhitha Reddy

\vspace{1pt}

\textit{Amrita School of Computing}\\
\textit{Amrita Vishwa Vidyapeetham}\\
Coimbatore, India\\
cb.sc.u4cse23206@cb.students.amrita.edu
\end{minipage}
&
\begin{minipage}[t]{0.30\textwidth}
\centering
\fontsize{8}{9}\selectfont
\textsuperscript{3rd} Shaik Abdus Sattar

\vspace{1pt}

\textit{Amrita School of Computing}\\
\textit{Amrita Vishwa Vidyapeetham}\\
Coimbatore, India\\
cb.sc.u4cse23245@cb.students.amrita.edu
\end{minipage}

\\[-1mm]

% VERTICAL GAP (MATCHING sample_paper.tex)
\multicolumn{3}{c}{\rule{0pt}{8mm}}\\[-1mm]

% ROW 2: AUTHORS 4, 5 (SYMMETRICALLY CENTERED)
\multicolumn{3}{c}{%
\begin{tabular}{@{}cc@{}}
\begin{minipage}[t]{0.36\textwidth}
\centering
\fontsize{8}{9}\selectfont
\textsuperscript{4th} Geethika G Nair

\vspace{1pt}

\textit{Amrita School of Computing}\\
\textit{Amrita Vishwa Vidyapeetham}\\
Coimbatore, India\\
cb.sc.u4cse23616@cb.students.amrita.edu
\end{minipage}
&
\begin{minipage}[t]{0.36\textwidth}
\centering
\fontsize{8}{9}\selectfont
\textsuperscript{5th} Meera J Vaishnav

\vspace{1pt}

\textit{Amrita School of Computing}\\
\textit{Amrita Vishwa Vidyapeetham}\\
Coimbatore, India\\
cb.sc.u4cse23633@cb.students.amrita.edu
\end{minipage}
\end{tabular}%
}

\end{tabular}

\endgroup
}

\maketitle

% =========================================================
% ABSTRACT
% =========================================================

\begin{abstract}
Graph matching and 3D geometric point-cloud registration are foundational challenges across computer vision, robotics, structural biology, and spatial network analysis. Evaluating structural affinity between complex relational systems requires identifying the optimal node correspondence via Weighted Common Subgraph (WCS) matching. Crucially, in contrast to opaque black-box Graph Neural Network (GNN) embeddings or compressed topological kernels that discard node-level traceability, recovering an explicit discrete bijective mapping $\mathbf{X}^* \in \mathcal{P}_L$ provides transparent, ground-truth explainability. In 3D geometric domains, this discrete 2D bijection directly provides the exact correspondence pairs $\mathcal{M} = \{(i_l, j_l)\}_{l=1}^L$ that unlock closed-form $\mathcal{O}(L)$ SVD Kabsch superposition (with $\mathcal{O}(1)$ $3 \times 3$ SVD), computing the optimal rigid transformation $(\mathbf{R}^*, \mathbf{t}^*)$ and coordinate RMSD in sub-millisecond time while bypassing iterative correspondence guessing (ICP) and soft fractional transport (POT, FGW). However, WCS matching is NP-hard. State-of-the-art non-convex relaxation schemes, such as the Graduated Non-Convexity and Concavity Procedure (GNCCP), suffer from severe computational bottlenecks in practice: evaluating the structural Hadamard objective without analytical gradients requires numerical finite differences scaling as $\Theta(N^5)$, while internal co-matching indicator tensor evaluations consume cubic work $\Theta(MN^2 + M^2N)$ per Frank-Wolfe iteration. Furthermore, prior polynomial relaxations (RLX1, RLX2) distort the continuous interior of the partial permutation polytope, inducing spurious stationary points and high asymptotic complexity. In this paper, we propose Fast-WCS, an analytical hierarchical 2D-to-3D optimization framework for attributed graph matching and closed-form physical superposition. Our framework introduces key analytical formulations and engineered algorithmic components: (1) a Rank-One Factorization (Lemma~1) of the co-matching indicator matrix $\mathbf{U} = \mathbf{r}\mathbf{r}^T$ that reduces internal tensor evaluations from cubic to quadratic $\mathcal{O}(MN + M^2)$ work; (2) an exact closed-form analytical structural gradient (Theorem~1) that eliminates finite-difference approximations and evaluates directly in $\Theta(N^3)$; (3) a Frank-Wolfe duality gap stopping test that adaptively certifies stationarity; (4) an exact quartic line search whose coefficients reuse cached matrix products and whose scalar cubic roots are solved in $\mathcal{O}(1)$ via Cardano's formula; (5) an exact Burkard Augmented Hungarian Oracle that guarantees exact cardinality $L$ in $\mathcal{O}((M+N-L)^3)$ time without heuristic truncation; and (6) a vectorized Level-3 BLAS combinatorial 2-opt local refinement that scores target replacement, target exchange, and source replacement moves simultaneously, breaking discrete mirror reflection ambiguities. We rigorously replicate the foundational synthetic and real-world image benchmarks of Yang et al. (IEEE TNNLS 2018), demonstrating that Fast-WCS achieves strictly cubic empirical scaling $\mathbf{\Theta(N^{3.0})}$ compared to $\mathcal{O}(N^{5.4})$ for prior GNCCP relaxations and $\mathcal{O}(N^{4.3})$ for classical spectral methods, running $19.3\times$ faster than RLX1 and scaling up to $M = 60$ in sub-$25\,\text{ms}$ with up to $97.1\%$ matching accuracy. On an extensive 28-configuration $(N, L)$ grid sweep spanning $N \in [20, 150]$ and $L/N \in [0.4, 1.0]$, Fast-WCS v3 executes in $52.2\,\text{ms}$ at $N=100$ ($10.5\times$ faster than RLX1F) and $206.7\,\text{ms}$ at $N=150$ ($165.4\times$ faster than RLX1F). On physical 3D geometric point-cloud superposition benchmarks across 100 multi-seed trials evaluated with untrimmed full-atom RMSD, Fast-WCS achieves $41.6\%\text{--}44.0\%$ accuracy with $\mathbf{37.0\%\text{--}42.0\%}$ sub-angstrom docking success ($< 1.0\,\text{\AA}$) and $26.0\%\text{--}27.0\%$ ultra-precision ($< 0.1\,\text{\AA}$) in $7.4\text{--}8.3\,\text{ms}$, whereas optimal transport and quadratic assignment baselines (POT, FGW, SciPy FAQ) completely fail ($0.0\text{--}1.0\%$ sub-angstrom recovery, median RMSD $> 2.7\,\text{\AA}$). Finally, on large-scale graphs up to $N = 100$, Fast-WCS evaluates in $140.0\,\text{ms}$ compared to $800.5\,\text{s}$ for numerical finite differences ($\mathbf{5,717.9\times}$ acceleration), establishing an efficient, mathematically rigorous foundation for hierarchical 2D graph matching and 3D physical superposition.
\end{abstract}

\begin{IEEEkeywords}
Attributed Graph Matching, Weighted Common Subgraph, GNCCP, Frank-Wolfe Optimization, Hierarchical 2D-to-3D Alignment, Bijective Correspondence, Closed-Form Kabsch Superposition, Point-Cloud Registration, Exact Quartic Line Search, Burkard Hungarian Oracle, Vectorized 2-Opt Refinement.
\end{IEEEkeywords}

% =========================================================
% I. INTRODUCTION
% =========================================================

\section{Introduction}
\label{sec:intro}

Quantifying structural similarity and discovering shared topological motifs between relational networks is a core mathematical challenge across modern scientific computing, computer vision, robotics, and structural biology \cite{conte2004, emmert2007}. In 3D geometric point-cloud registration, visual landmark tracking, and physical macromolecular modeling, evaluating similarity between two complex systems requires finding maximum common substructures while accommodating flexible deformations, rotational invariance, and varying background noise \cite{zaslavskiy2009, vishwanathan2010, leordeanu2005}. Within this broad domain, Weighted Common Subgraph (WCS) matching provides a principled mathematical formulation for establishing discrete correspondences across attributed relational structures.

\begin{figure*}[t]
\centering
\includegraphics[width=0.98\textwidth]{fig_workflow_pipeline.pdf}
\caption{End-to-end architecture of the proposed Fast-WCS hierarchical optimization framework. Phase~1 extracts attributed graphs $G$ and $H$ from relational data or 3D coordinate point clouds, processes them through the analytical Fast-WCS homotopy engine (leveraging Lemma~1 Rank-1 Factorization, Theorem~1 exact closed-form gradient, Hungarian LMO, and a Frank-Wolfe gap stopping test), and computes an exact discrete bijective node correspondence $\mathbf{X}^* \in \mathcal{P}_L$. Phase~2 exploits this interpretable mapping to execute a closed-form $\mathcal{O}(L)$ SVD Kabsch superposition of physical 3D coordinates, computing the optimal rigid transformation $(\mathbf{R}^*, \mathbf{t}^*)$ and geometric RMSD in sub-millisecond time. The resulting 2D topological and 3D spatial similarity metrics provide certified geometric registration without iterative correspondence guessing.}
\label{fig:workflow_pipeline}
\end{figure*}

Recently, deep Graph Neural Networks (GNNs) have emerged as popular tools for relational learning. However, GNNs operate as opaque black boxes that compress intricate graph topologies into low-dimensional continuous vector embeddings. This compression discards node-level isomorphism guarantees, suffers from catastrophic oversmoothing and oversquashing on larger graphs, and crucially fails to produce an auditable correspondence between individual vertices. In safety-critical applications such as robotic manipulation, medical imaging, or molecular docking, scientists and engineers cannot inspect \textit{why} two graphs or shapes are declared similar by a black-box embedding. In contrast, classical combinatorial matching establishes an explicit discrete bijective mapping $\mathbf{X}^* \in \mathcal{P}_L$ between a subset of $L$ nodes in graph $G$ and $L$ nodes in graph $H$. This correspondence provides total, ground-truth \textbf{explainability}: each node assignment is transparently verifiable. 

Furthermore, in physical domains where vertices possess embedded 3D Cartesian coordinates (such as articulated point clouds, robotic visual landmarks, or atomic coordinates), an exact discrete bijection $\mathbf{X}^*$ directly unlocks a closed-form $\mathcal{O}(L)$ Singular Value Decomposition (SVD) Kabsch superposition. Traditional 3D registration frameworks, such as the Iterative Closest Point (ICP) algorithm, alternate between heuristic spatial nearest-neighbor assignment and rigid pose estimation, frequently becoming trapped in poor local minima when rotations are large or outliers exist. Similarly, modern optimal transport formulations (e.g., Partial Optimal Transport or Fused Gromov-Wasserstein) produce continuous fractional coupling matrices $\mathbf{T} \in [0, 1]^{M \times N}$, which blur correspondences and lack discrete rigidity. In stark contrast, our hierarchical formulation decouples relational correspondence from spatial pose: Phase~1 recovers an exact discrete bijection $\mathbf{X}^* \in \mathcal{P}_L$ based on pairwise topological metrics and vertex attributes, and Phase~2 evaluates the optimal rigid rotation $\mathbf{R}^* \in \mathrm{SO}(3)$, translation $\mathbf{t}^* \in \mathbb{R}^3$, and coordinate Root-Mean-Square Deviation (RMSD) in closed form in microseconds!

Within this general mathematical formulation, the problem of quantifying structural affinity between two attributed weighted graphs $G$ and $H$ containing $M$ and $N$ nodes ($M \leq N$) reduces to the Weighted Common Subgraph (WCS) matching problem for a prescribed subgraph size $L \leq \min(M, N)$ \cite{yang2018algorithm}. The resulting optimal assignment matrix $\mathbf{X}^* \in \mathcal{P}_L$ yields both an objective distance metric and an explicit bijective correspondence mapping matched nodes between the systems, while naturally isolating $M - L$ nodes in $G$ and $N - L$ nodes in $H$ as unassigned outliers.

However, WCS matching is an NP-hard combinatorial problem that generalizes the maximum common subgraph isomorphism and quadratic assignment problems (QAP) \cite{garey1979, loiola2007}. Classical approaches to subgraph matching frequently decouple the optimization into two independent stages: solving a full graph matching problem between $G$ and $H$, followed by heuristic vertex pruning to retain the $L$ highest-scoring vertices \cite{zaslavskiy2009, gold1996}. Unfortunately, optimizing global alignment does not imply optimality over the embedded subgraphs, frequently discarding high-affinity functional subgraphs in the presence of unassigned outlier nodes and background noise.

To resolve this limitation, Liu and Qiao \cite{liu2015graduated} and Yang et al. \cite{yang2018algorithm} formulated WCS matching directly on the partial permutation polytope $\mathcal{P}_L$. By relaxing the discrete binary assignment constraints into their continuous convex hull $\mathcal{D}_L$ and introducing an annealing continuation parameter $\zeta \in [-1, 1]$, the Graduated Non-Convexity and Concavity Procedure (GNCCP) deforms a trivially solvable strictly convex function at $\zeta = +1$ into the original NP-hard non-convex objective at $\zeta = 0$, terminating at a concave problem at $\zeta = -1$ whose extreme points are provably discrete permutation matrices \cite{yang2018algorithm}.

Despite its theoretical elegance, existing implementations of the GNCCP framework suffer from severe computational and analytical bottlenecks:
\begin{enumerate}
    \item \textbf{Analytical Gradient Intractability vs. Distortion:} The structural mismatch term in the WCS objective couples the assignment matrix $\mathbf{X}$ with an internal co-matching indicator matrix $\mathbf{U} = \mathbf{X}\mathbf{1}_{N\times N}\mathbf{X}^T$ via an entry-wise Hadamard product. In foundational implementations \cite{yang2018algorithm}, calculating the gradient of this Hadamard term directly was considered analytically intractable. To circumvent this, Yang et al. \cite{yang2018algorithm} proposed polynomial relaxations (RLX1 and RLX2); however, their higher-order expansions distort the continuous landscape over the interior $\mathcal{D}_L \setminus \mathcal{P}_L$, introducing artificial stationary points and requiring expensive inner solvers scaling as $\mathcal{O}(N^{5.4})$. Alternatively, evaluating the true objective gradient numerically via finite differences requires $\Theta(MN)$ full objective calls per step, causing per-iteration complexity to explode to $\Theta(N^5)$.
    \item \textbf{Cubic Tensor Co-Matching Work:} The internal matrix $\mathbf{U}$, which tracks whether pairs of nodes $(i, k)$ in graph $G$ are simultaneously selected in the matching, naively involves multiplying an $M \times N$ matrix by an $N \times N$ all-ones matrix and an $N \times M$ matrix. This operation consumes $\Theta(MN^2 + M^2N)$ arithmetic operations at every gradient step.
    \item \textbf{Inexact Line Searches and Heuristic Truncations:} Prior implementations rely on expensive numerical grid evaluations or heuristic Hungarian column dropping (RLX1F) that violate mathematical cardinality constraints and degrade matching accuracy under negative gradient terms.
    \item \textbf{Combinatorial Post-Hoc Bottlenecks:} Transitioning from continuous relaxations to discrete physical bijections previously required either unguided Hungarian projection or slow $\mathcal{O}(L^3)$ nested loops for swap refinements, creating major runtime bottlenecks and failing to resolve mirror reflection ambiguities in 3D.
\end{enumerate}

\subsection{Key Contributions of This Work}

In this paper, we propose \textbf{Fast-WCS}, an analytical hierarchical 2D-to-3D optimization framework for general attributed weighted graphs and 3D physical point clouds. It evaluates the original WCS objective and enforces the prescribed partial-assignment cardinality while using an exact structural gradient, an exact linear assignment oracle, and closed-form 3D superposition. As illustrated in Fig.~\ref{fig:workflow_pipeline}, our methodological contributions are:

\begin{itemize}
    \item \textbf{Lemma~1 (Rank-One Factorization):} We prove that the internal co-matching matrix $\mathbf{U} = \mathbf{X}\mathbf{1}_{N\times N}\mathbf{X}^T$ factors into an outer product of the row-sum vector $\mathbf{U} = \mathbf{r}\mathbf{r}^T$, where $\mathbf{r} = \mathbf{X}\mathbf{1}_N$. This drops the evaluation cost from cubic $\Theta(MN^2 + M^2N)$ to quadratic $\mathcal{O}(MN + M^2)$ work, eliminating an entire cubic matrix product per step.
    \item \textbf{Theorem~1 (Closed-Form Analytical Gradient):} We derive the exact analytical gradient $\nabla_{\mathbf{X}} f(\mathbf{X})$ of the true unrelaxed structural mismatch objective in closed form. By combining trace differentials with Hadamard properties and Lemma~1, the gradient evaluates in strictly $\Theta(N^3)$ operations without any finite-difference approximations or polynomial relaxation distortion.
    \item \textbf{Frank-Wolfe Duality Gap Stopping:} At each homotopy level, the implementation computes the Frank-Wolfe duality gap and stops the inner loop adaptively when its absolute value is below the configured tolerance, eliminating redundant iterations.
    \item \textbf{Theorem~3 (Exact Quartic Line Search):} Along a Frank-Wolfe segment, $J_\zeta(\lambda)$ is an exact quartic polynomial. The implementation forms its coefficients using products cached during gradient evaluation and computes its real stationary points in $\mathcal{O}(1)$ closed form via Cardano's formula with a numerical fallback, replacing all grid search approximations.
    \item \textbf{Exact Burkard Augmented Hungarian Oracle:} We formulate an exact $(M+N-L) \times (M+N-L)$ augmented cost matrix with an infinite penalty slack block, rigorously guaranteeing exact cardinality $L$ in $\mathcal{O}((M+N-L)^3)$ time without requiring lossy column dropping or ad-hoc positivity shifts.
    \item \textbf{Vectorized Combinatorial Local Refinement:} After Hungarian discretization, the implementation scores target replacements, pairwise target exchanges, and source replacements in vectorized Level-3 BLAS matrix operations, actively resolving discrete mirror reflection ambiguities ($\det(\mathbf{R}) = -1 \to +1$).
    \item \textbf{Hierarchical 2D Bijection to 3D Kabsch Superposition:} We establish a principled two-phase hierarchy where the discrete bijection $\mathbf{X}^* \in \mathcal{P}_L$ recovered in Phase~1 provides the exact correspondences that unlock closed-form $\mathcal{O}(L)$ SVD Kabsch superposition in Phase~2, eliminating heuristic correspondence guessing.
    \item \textbf{Extensive Benchmarking Across 2D and 3D Regimes:} We validate Fast-WCS across synthetic graph benchmarks, real-world image landmarks, an extensive 28-configuration $(N, L)$ grid sweep spanning $N \in [20, 150]$, and 100 multi-seed 3D point-cloud trials with untrimmed full-atom RMSD. Fast-WCS achieves strictly cubic scaling $\mathbf{\Theta(N^{3.0})}$, runs up to $165.4\times$ faster than RLX1F at $N=150$ ($5,717.9\times$ faster than finite differences at $N=100$), and attains $\mathbf{37.0\%\text{--}42.0\%}$ sub-angstrom precision ($< 1.0\,\text{\AA}$) in $7.4\text{--}8.3\,\text{ms}$, where modern optimal transport and QAP baselines (POT, FGW, SciPy FAQ) achieve $0.0\%$ sub-angstrom recovery.
\end{itemize}

% =========================================================
% II. RELATED WORK AND FOUNDATIONS
% =========================================================

\section{Related Work and Theoretical Foundations}
\label{sec:related_work}

Graph matching has a rich history spanning computer vision, pattern recognition, and chemoinformatics \cite{conte2004, emmert2007}. Mathematically, matching two graphs is formulated as an instance of the Quadratic Assignment Problem (QAP) in Koopmans-Beckmann \cite{koopmans1957} or Lawler form \cite{lawler1963}.

\subsection{Quadratic Assignment and Affinity Matrix Formulations}

Given two graphs $G = (V_G, E_G)$ and $H = (V_H, E_H)$ with adjacency matrices $\mathbf{A}_G \in \mathbb{R}^{M \times M}$ and $\mathbf{A}_H \in \mathbb{R}^{N \times N}$, the Koopmans-Beckmann QAP seeks an assignment $\mathbf{X}$ maximizing:
\begin{equation}
\max_{\mathbf{X} \in \Pi} \operatorname{tr}(\mathbf{X}^T \mathbf{A}_G \mathbf{X} \mathbf{A}_H),
\label{eq:koopmans}
\end{equation}
where $\Pi$ denotes the set of permutation matrices. In the more general Lawler formulation adopted by computer vision methods, an affinity matrix $\mathbf{K} \in \mathbb{R}^{MN \times MN}$ is constructed, where diagonal entries $K_{ia, ia}$ encode vertex similarity and off-diagonal entries $K_{ia, jb}$ encode pairwise edge compatibility between edge $(i, j) \in E_G$ and edge $(a, b) \in E_H$ \cite{leordeanu2005, gold1996}. The matching problem then takes the form:
\begin{equation}
\max_{\mathbf{x} \in \{0, 1\}^{MN}} \mathbf{x}^T \mathbf{K} \mathbf{x} \quad \text{s.t.} \quad \sum_{a=1}^N x_{ia} \leq 1, \; \sum_{i=1}^M x_{ia} \leq 1,
\label{eq:lawler}
\end{equation}
where $\mathbf{x} = \operatorname{vec}(\mathbf{X})$.

While the Lawler formulation is expressive, its storage complexity is $\mathcal{O}(M^2 N^2)$, requiring $\mathcal{O}(N^4)$ memory and $\mathcal{O}(N^4)$ matrix-vector multiplications per iteration \cite{zaslavskiy2009}. For example, comparing two moderately sized drug compounds with $N = 50$ atoms requires storing an affinity matrix of dimension $2,500 \times 2,500$ ($6.25 \times 10^6$ entries), consuming gigabytes of memory and rendering pairwise kernel evaluations intractable.

Classical algorithms addressing \eqref{eq:lawler} include the Spectral Matching (SM) algorithm of Leordeanu and Hebert \cite{leordeanu2005}, which relaxes $\mathbf{x}$ to the unit sphere and computes the leading eigenvector of $\mathbf{K}$, and the Graduated Assignment (GA) algorithm of Gold and Rangarajan \cite{gold1996}, which incorporates deterministic annealing with Softmax normalizations. Both methods suffer from $\mathcal{O}(N^4)$ affinity matrix overhead and struggle in the presence of unassigned outlier nodes.

\subsection{Path-Following and Convex-Concave Procedures}

To avoid the $\mathcal{O}(N^4)$ affinity matrix bottleneck, Zaslavskiy et al. \cite{zaslavskiy2009} introduced the Extended Path Following (EPF) algorithm for adjacency-based graph matching, operating directly on $\mathbf{A}_G$ and $\mathbf{A}_H$ in $\mathcal{O}(N^2)$ space. EPF constructs a convex relaxation $J_{\text{cvx}}(\mathbf{X})$ and a concave relaxation $J_{\text{ccv}}(\mathbf{X})$ and tracks the solution along the homotopy path $(1 - \alpha) J_{\text{cvx}} + \alpha J_{\text{ccv}}$ as $\alpha$ increases from $0$ to $1$.

Building on this principle, Liu and Qiao \cite{liu2015graduated} established the Graduated Non-Convexity and Concavity Procedure (GNCCP) for full graph matching, demonstrating that any quadratic objective $\mathbf{x}^T \mathbf{K} \mathbf{x}$ can be embedded into a continuous homotopy path with guaranteed convergence to discrete extreme points. Yang et al. \cite{yang2018algorithm} subsequently extended GNCCP to the Weighted Common Subgraph (WCS) matching problem for arbitrary subgraph sizes $L \leq \min(M, N)$.

\begin{figure*}[t]
\centering
\includegraphics[width=0.78\textwidth]{fig_wcs_concept.pdf}
\caption{Foundational formulation of Weighted Common Subgraph (WCS) matching and algebraic subgraph extraction (replicating the formulation of Yang et al.~\cite{yang2018algorithm}). Top: Establishing a discrete bijection between $L = 4$ vertices in $G$ ($M = 5$) and $H$ ($N = 6$) via the partial permutation matrix $\mathbf{X} \in \mathcal{P}_L$. Bottom (a--d): Algebraic definition of the co-matching matrix $\mathbf{U} = \mathbf{X}\mathbf{1}_{N\times N}\mathbf{X}^T$, subgraph extraction $\mathbf{A}_{\text{sub}G} = \mathbf{U} \circ \mathbf{A}_G$, aligned target subgraph $\mathbf{A}_{\text{sub}H} = \mathbf{X}\mathbf{A}_H\mathbf{X}^T$, and quadratic structural mismatch $H_0(\mathbf{X}) = \|\mathbf{A}_{\text{sub}G} - \mathbf{A}_{\text{sub}H}\|_F^2$.}
\label{fig:wcs_concept}
\end{figure*}

\subsection{Yang et al. (2018) Polynomial Relaxations and Limitations}

In formulating WCS matching, Yang et al. \cite{yang2018algorithm} recognized that the exact structural mismatch term involves a Hadamard product:
\begin{equation}
H_0(\mathbf{X}) = \big\| \mathbf{U} \circ \mathbf{A}_G - \mathbf{X}\mathbf{A}_H\mathbf{X}^T \big\|_F^2,
\label{eq:yang_h0}
\end{equation}
where $\mathbf{U} = \mathbf{X}\mathbf{1}_{N\times N}\mathbf{X}^T$. Crucially, Yang et al. asserted that ``in implementation, it is difficult to directly calculate $\nabla H_0(\mathbf{X})$, which involves a Hadamard product.'' Because they could not differentiate $H_0(\mathbf{X})$ analytically, they replaced $H_0(\mathbf{X})$ with polynomial relaxations:

\begin{enumerate}
    \item \textbf{RLX1 Formulation:} By leveraging algebraic properties that hold exclusively on binary matrices $\mathbf{X} \in \mathcal{P}_L$ (namely $\mathbf{U}\circ\mathbf{U} = \mathbf{U}$ and $\mathbf{U} \circ (\mathbf{X}\mathbf{A}_H\mathbf{X}^T) = \mathbf{X}\mathbf{A}_H\mathbf{X}^T$), they expanded $H_0(\mathbf{X})$ into:
    \begin{align}
    H_1(\mathbf{X}) = &\operatorname{tr}((\mathbf{A}_G \circ \mathbf{A}_G)\mathbf{U}^T) - 2\operatorname{tr}(\mathbf{A}_G\mathbf{X}\mathbf{A}_H^T\mathbf{X}^T) \notag\\
    &+ \operatorname{tr}(\mathbf{X}\mathbf{A}_H\mathbf{X}^T\mathbf{X}\mathbf{A}_H^T\mathbf{X}^T),
    \label{eq:rlx1}
    \end{align}
    with gradient:
    \begin{align}
    \nabla H_1(\mathbf{X}) = &(\mathbf{A}_G^T \circ \mathbf{A}_G^T + \mathbf{A}_G \circ \mathbf{A}_G)\mathbf{X}\mathbf{1}_{N\times N} \notag\\
    &- 2(\mathbf{A}_G^T\mathbf{X}\mathbf{A}_H + \mathbf{A}_G\mathbf{X}\mathbf{A}_H^T) \notag\\
    &+ 2(\mathbf{X}\mathbf{A}_H\mathbf{X}^T\mathbf{X}\mathbf{A}_H^T + \mathbf{X}\mathbf{A}_H^T\mathbf{X}^T\mathbf{X}\mathbf{A}_H).
    \label{eq:rlx1_grad}
    \end{align}
    \item \textbf{RLX2 Formulation:} They substituted $\mathbf{U}\circ\mathbf{A}_G$ with $(\mathbf{X}\mathbf{X}^T)\mathbf{A}_G(\mathbf{X}\mathbf{X}^T)$, expanding into a sixth-order polynomial $H_2(\mathbf{X})$ with nine distinct matrix products.
    \item \textbf{RLX1F Truncation Heuristic:} When $L < M \leq N$, the partial assignment linear minimization oracle cannot be solved by the standard rectangular Hungarian algorithm. Yang et al. therefore solved an unconstrained $M$-assignment and dropped the $M - L$ entries with smallest $-\nabla J_\zeta$.
\end{enumerate}

These approximations have severe drawbacks: on the continuous interior $\mathcal{D}_L \setminus \mathcal{P}_L$, $H_1(\mathbf{X}) \neq H_0(\mathbf{X})$ and $H_2(\mathbf{X}) \neq H_0(\mathbf{X})$. The fourth-order terms in $\nabla H_1(\mathbf{X})$ and sixth-order terms in $\nabla H_2(\mathbf{X})$ distort the descent trajectory, create spurious local minima, and force inner solvers to scale as $\mathcal{O}(N^{5.4})$. Furthermore, the RLX1F heuristic truncation causes substantial accuracy degradation.

Fast-WCS fundamentally resolves these challenges: by proving Lemma~1 and Theorem~1, we calculate the exact closed-form gradient of $H_0(\mathbf{X})$ directly in $\Theta(N^3)$, eliminating all polynomial relaxations and heuristic truncations.

% =========================================================
% III. PROBLEM FORMULATION
% =========================================================

\section{Problem Formulation}
\label{sec:formulation}

Let $G = (V_G, E_G, \mathbf{A}_G, \mathbf{F}_G)$ and $H = (V_H, E_H, \mathbf{A}_H, \mathbf{F}_H)$ represent two general attributed weighted graphs, with vertex sets $|V_G| = M$ and $|V_H| = N$ ($M \leq N$).

\subsection{Weighted Attributed Graph Representation}

The relational topology of graphs $G$ and $H$ is encoded by non-negative symmetric weighted adjacency matrices $\mathbf{A}_G \in \mathbb{R}_{\geq 0}^{M \times M}$ and $\mathbf{A}_H \in \mathbb{R}_{\geq 0}^{N \times N}$. For any pair of vertices $(i, j)$, the edge weight $A_{ij} \geq 0$ quantifies connection strength, interaction affinity, or inverse spatial distance between nodes, with $A_{ij} = 0$ denoting the absence of an edge.

Vertex-level semantic properties are captured by attribute matrices $\mathbf{F}_G \in \mathbb{R}^{M \times d}$ and $\mathbf{F}_H \in \mathbb{R}^{N \times d}$, where each row represents a $d$-dimensional normalized feature vector encoding domain-specific node properties (e.g., chemical element types, hybridization states, and formal charges in molecular graphs, or 2D/3D spatial coordinates in geometric point clouds). A unary discrepancy cost matrix $\mathbf{C} \in \mathbb{R}_{\geq 0}^{M \times N}$ is defined by pairwise Euclidean feature distances:
\begin{equation}
C_{ij} = \|\mathbf{F}_{G, i} - \mathbf{F}_{H, j}\|_2,
\label{eq:unary_cost}
\end{equation}
quantifying the penalty for placing vertex $i \in V_G$ in correspondence with vertex $j \in V_H$.

\subsection{Partial Permutation Polytope}

Let $L \in \{1, 2, \ldots, M\}$ denote the target size of the common subgraph to be identified. The discrete correspondence is parameterized by a binary assignment matrix $\mathbf{X} \in \{0, 1\}^{M \times N}$, where $X_{ij} = 1$ indicates that vertex $i \in V_G$ corresponds to vertex $j \in V_H$.

\begin{definition}[Partial Permutation Polytope $\mathcal{P}_L$]
The discrete partial permutation polytope $\mathcal{P}_L$ is defined as:
\begin{align}
\mathcal{P}_L = \Big\{ \mathbf{X} \in \{0, 1\}^{M \times N} \;\Big|\; 
&\sum_{j=1}^N X_{ij} \leq 1, \; \forall i \in \{1, \ldots, M\}, \notag\\
&\sum_{i=1}^M X_{ij} \leq 1, \; \forall j \in \{1, \ldots, N\}, \notag\\
&\sum_{i=1}^M \sum_{j=1}^N X_{ij} = L \Big\}.
\label{eq:polytope_discrete}
\end{align}
\end{definition}

In vector notation, the constraints are $\mathbf{X}\mathbf{1}_N \leq \mathbf{1}_M$, $\mathbf{X}^T\mathbf{1}_M \leq \mathbf{1}_N$, and $\mathbf{1}_M^T\mathbf{X}\mathbf{1}_N = L$, where $\mathbf{1}_K$ denotes the all-ones vector in $\mathbb{R}^K$. As illustrated in Fig.~\ref{fig:wcs_concept}, any $\mathbf{X} \in \mathcal{P}_L$ selects exactly $L$ vertices in $G$ and $L$ vertices in $H$, leaving $M - L$ vertices in $G$ and $N - L$ vertices in $H$ unassigned.

The continuous relaxation $\mathcal{D}_L$ replaces the discrete set $\{0, 1\}^{M \times N}$ with the unit hypercube $[0, 1]^{M \times N}$:
\begin{align}
\mathcal{D}_L = \Big\{ \mathbf{X} \in [0, 1]^{M \times N} \;\Big|\; 
&\mathbf{X}\mathbf{1}_N \leq \mathbf{1}_M, \; \mathbf{X}^T\mathbf{1}_M \leq \mathbf{1}_N, \notag\\
&\mathbf{1}_M^T\mathbf{X}\mathbf{1}_N = L \Big\}.
\label{eq:polytope_continuous}
\end{align}
By the Birkhoff-von Neumann theorem extended to rectangular doubly sub-stochastic matrices, $\mathcal{D}_L$ is a compact, convex polytope whose extreme points (vertices) are precisely the elements of $\mathcal{P}_L$ \cite{birkhoff1946}.

\begin{figure*}[t]
\centering
\includegraphics[width=0.95\textwidth]{fig_wcs_objective.pdf}
\caption{Analytical structure of the Fast-WCS objective, Lemma~1 Rank-1 Factorization, and Theorem~1 Exact Analytical Gradient. (a)~Conventional naive evaluation of the co-matching indicator matrix $\mathbf{U} = \mathbf{X}\mathbf{1}_{N\times N}\mathbf{X}^T$ requires multiplying by an $N \times N$ dense all-ones matrix, consuming $\Theta(MN^2 + M^2N)$ operations. (b)~Proposed Lemma~1 Rank-1 Factorization computes the row-sum vector $\mathbf{r} = \mathbf{X}\mathbf{1}_N$ and forms $\mathbf{U} = \mathbf{r}\mathbf{r}^T$ as an outer product in strictly $\mathcal{O}(MN + M^2)$ work, eliminating intermediate cubic tensor contractions. (c)~Hadamard structural extraction $\mathbf{A}_{\text{sub}G} = \mathbf{U} \circ \mathbf{A}_G$ versus target projection $\mathbf{A}_{\text{sub}H} = \mathbf{X}\hat{\mathbf{A}}_H\mathbf{X}^T$. (d)~Theorem~1 exact closed-form gradient $\nabla_{\mathbf{X}} F(\mathbf{X}) = 4(\mathbf{R} \circ \hat{\mathbf{A}}_G)\mathbf{r}\mathbf{1}_N^T - 4\mathbf{R}\mathbf{X}\hat{\mathbf{A}}_H$ evaluating in strictly $\Theta(N^3)$ operations via three dense matrix multiplications with zero finite-difference approximations.}
\label{fig:wcs_objective}
\end{figure*}

\subsection{Weighted Common Subgraph Objective}

The Weighted Common Subgraph optimization seeks an assignment matrix $\mathbf{X} \in \mathcal{P}_L$ minimizing the combined topological mismatch and unary attribute cost:
\begin{equation}
F(\mathbf{X}) = \alpha\, f(\mathbf{X}) + (1 - \alpha)\operatorname{tr}(\mathbf{C}^T\mathbf{X}),
\label{eq:wcs_full_objective}
\end{equation}
where $\alpha \in [0, 1]$ is a trade-off parameter balancing topological structure against node identity, $\operatorname{tr}(\mathbf{C}^T\mathbf{X}) = \sum_{i=1}^M \sum_{j=1}^N C_{ij} X_{ij}$ measures unary attribute discrepancy, and $f(\mathbf{X})$ is the structural mismatch function:
\begin{equation}
f(\mathbf{X}) = \big\| \mathbf{U} \circ \hat{\mathbf{A}}_G - \mathbf{X}\hat{\mathbf{A}}_H\mathbf{X}^T \big\|_F^2.
\label{eq:f_structural}
\end{equation}
Here, $\hat{\mathbf{A}}_G = \mathbf{A}_G / \|\mathbf{A}_G\|_F$ and $\hat{\mathbf{A}}_H = \mathbf{A}_H / \|\mathbf{A}_H\|_F$ are Frobenius-normalized adjacency matrices, $\circ$ denotes the Hadamard (entry-wise) product, and $\mathbf{U} \in \mathbb{R}^{M \times M}$ is the co-matching indicator matrix defined as:
\begin{equation}
\mathbf{U} = \mathbf{X}\mathbf{1}_{N\times N}\mathbf{X}^T.
\label{eq:U_naive}
\end{equation}

For any binary matrix $\mathbf{X} \in \mathcal{P}_L$, entry $U_{ik} = (\sum_{j=1}^N X_{ij})(\sum_{l=1}^N X_{kl})$ equals $1$ if and only if vertex $i$ and vertex $k$ in $G$ are both assigned to vertices in $H$, and $0$ otherwise. Thus, the Hadamard product $\mathbf{U} \circ \hat{\mathbf{A}}_G$ extracts the induced subgraph of $G$ restricted to the matched vertices, while $\mathbf{X}\hat{\mathbf{A}}_H\mathbf{X}^T$ projects the adjacency of $H$ onto the indexed coordinates of $G$, as illustrated in Fig.~\ref{fig:wcs_objective}c.

% =========================================================
% IV. ANALYTICAL FOUNDATIONS AND THEORETICAL PROOFS
% =========================================================

\section{Analytical Foundations and Theoretical Proofs}
\label{sec:theory}

In this section, we present the formal mathematical derivations and complete proofs that form the theoretical foundation of Fast-WCS.

\subsection{Rank-One Factorization of Co-Matching Matrix}

Evaluating $\mathbf{U}$ via \eqref{eq:U_naive} naively requires multiplying an $M \times N$ matrix by an $N \times N$ matrix followed by an $N \times M$ matrix, consuming $\Theta(MN^2 + M^2N)$ arithmetic operations per step.

\begin{lemma}[Rank-One Factorization of $\mathbf{U}$]
\label{lem:rank1}
Let $\mathbf{r} = \mathbf{X}\mathbf{1}_N \in \mathbb{R}^M$ denote the row-sum vector of the assignment matrix $\mathbf{X}$, where $r_i = \sum_{j=1}^N X_{ij}$. Then:
\begin{equation}
\mathbf{U} = \mathbf{X}\mathbf{1}_{N\times N}\mathbf{X}^T = \mathbf{r}\mathbf{r}^T.
\label{eq:lemma1_result}
\end{equation}
\end{lemma}

\begin{proof}
The dense all-ones matrix $\mathbf{1}_{N\times N} \in \mathbb{R}^{N\times N}$ factors into the outer product of two $N$-dimensional all-ones vectors:
\begin{equation}
\mathbf{1}_{N\times N} = \mathbf{1}_N \mathbf{1}_N^T.
\end{equation}
Substituting this factorized form into the definition of $\mathbf{U}$:
\begin{align}
\mathbf{U} &= \mathbf{X} \big(\mathbf{1}_N \mathbf{1}_N^T\big) \mathbf{X}^T \notag\\
&= \big(\mathbf{X}\mathbf{1}_N\big) \big(\mathbf{X}\mathbf{1}_N\big)^T \notag\\
&= \mathbf{r} \mathbf{r}^T.
\end{align}
Expanding in index notation, for each entry $(i, k) \in \{1, \ldots, M\}^2$:
\begin{align}
U_{ik} &= \sum_{j=1}^N \sum_{l=1}^N X_{ij} (\mathbf{1}_{N\times N})_{jl} X_{kl} \notag\\
&= \left(\sum_{j=1}^N X_{ij}\right)\!\left(\sum_{l=1}^N X_{kl}\right) = r_i r_k.
\end{align}
This completes the proof.
\end{proof}

\begin{remark}[Computational Complexity Savings]
Computing $\mathbf{r} = \mathbf{X}\mathbf{1}_N$ requires $MN$ additions. The outer product $\mathbf{r}\mathbf{r}^T$ requires $M^2$ multiplications. The total complexity is $\mathcal{O}(MN + M^2)$. For square graphs where $M = N = n$, this reduces the computation from $2n^3$ operations to $n^2 + n^2 = 2n^2$ operations, entirely eliminating cubic work from the co-matching step.
\end{remark}

\subsection{Closed-Form Analytical Structural Gradient}

We now derive the exact analytical gradient of the structural mismatch cost $f(\mathbf{X})$ without making any polynomial relaxations or finite-difference approximations.

\begin{theorem}[Closed-Form Analytical Structural Gradient]
\label{thm:gradient}
Let the structural residual matrix be defined as:
\begin{equation}
\mathbf{R} = (\mathbf{U} \circ \hat{\mathbf{A}}_G) - (\mathbf{X}\hat{\mathbf{A}}_H\mathbf{X}^T) \in \mathbb{R}^{M \times M}.
\label{eq:residual}
\end{equation}
Assume $\hat{\mathbf{A}}_G$ and $\hat{\mathbf{A}}_H$ are symmetric, as required by the v3 implementation.
The exact gradient of $f(\mathbf{X}) = \|\mathbf{R}\|_F^2$ with respect to $\mathbf{X}$ is given by:
\begin{equation}
\nabla_{\mathbf{X}} f(\mathbf{X}) = 4\big(\mathbf{R} \circ \hat{\mathbf{A}}_G\big)\mathbf{r}\mathbf{1}_N^T - 4\mathbf{R}\mathbf{X}\hat{\mathbf{A}}_H \in \mathbb{R}^{M \times N},
\label{eq:thm1_result}
\end{equation}
where $\mathbf{r} = \mathbf{X}\mathbf{1}_N$. Furthermore, the first term is a rank-one column broadcast matrix where row $i$ repeats the scalar value $4[(\mathbf{R} \circ \hat{\mathbf{A}}_G)\mathbf{r}]_i$ across all $N$ columns.
\end{theorem}

\begin{proof}
The structural cost is expressed using the matrix trace inner product as:
\begin{equation}
f(\mathbf{X}) = \operatorname{tr}(\mathbf{R}^T \mathbf{R}) = \operatorname{tr}(\mathbf{R}\mathbf{R}),
\end{equation}
since $\mathbf{R}$ is symmetric ($\mathbf{U}$, $\hat{\mathbf{A}}_G$, and $\mathbf{X}\hat{\mathbf{A}}_H\mathbf{X}^T$ are symmetric). Taking the total differential:
\begin{equation}
\mathrm{d}f = 2\operatorname{tr}(\mathbf{R}^T \mathrm{d}\mathbf{R}) = 2\operatorname{tr}(\mathbf{R} \, \mathrm{d}\mathbf{R}).
\label{eq:df_trace}
\end{equation}
From the definition of $\mathbf{R}$ in \eqref{eq:residual}:
\begin{equation}
\mathrm{d}\mathbf{R} = \mathrm{d}(\mathbf{U} \circ \hat{\mathbf{A}}_G) - \mathrm{d}(\mathbf{X}\hat{\mathbf{A}}_H\mathbf{X}^T).
\label{eq:dR}
\end{equation}
We evaluate the two differential terms separately.

\textit{Part 1: Hadamard differential $\mathrm{d}(\mathbf{U} \circ \hat{\mathbf{A}}_G)$.} Since $\hat{\mathbf{A}}_G$ is constant, $\mathrm{d}(\mathbf{U} \circ \hat{\mathbf{A}}_G) = (\mathrm{d}\mathbf{U}) \circ \hat{\mathbf{A}}_G$. Using Lemma~\ref{lem:rank1}, $\mathbf{U} = \mathbf{r}\mathbf{r}^T$ with $\mathbf{r} = \mathbf{X}\mathbf{1}_N$. By the product rule:
\begin{equation}
\mathrm{d}\mathbf{U} = (\mathrm{d}\mathbf{r})\mathbf{r}^T + \mathbf{r}(\mathrm{d}\mathbf{r})^T = (\mathrm{d}\mathbf{X}\mathbf{1}_N)\mathbf{r}^T + \mathbf{r}(\mathrm{d}\mathbf{X}\mathbf{1}_N)^T.
\end{equation}
Substituting into the Hadamard trace identity $\operatorname{tr}(\mathbf{A}(\mathbf{B}\circ\mathbf{C})) = \operatorname{tr}((\mathbf{A}\circ\mathbf{B})\mathbf{C})$:
\begin{align}
\operatorname{tr}\big(\mathbf{R} \, \mathrm{d}(\mathbf{U}\circ\hat{\mathbf{A}}_G)\big) &= \operatorname{tr}\big(\mathbf{R} \, ((\mathrm{d}\mathbf{U})\circ\hat{\mathbf{A}}_G)\big) \notag\\
&= \operatorname{tr}\big((\mathbf{R} \circ \hat{\mathbf{A}}_G) \, \mathrm{d}\mathbf{U}\big).
\end{align}
Let $\mathbf{S} = \mathbf{R} \circ \hat{\mathbf{A}}_G$. Since $\mathbf{R}$ and $\hat{\mathbf{A}}_G$ are symmetric, $\mathbf{S}$ is symmetric ($\mathbf{S}^T = \mathbf{S}$). Expanding with $\mathrm{d}\mathbf{U}$:
\begin{align}
\operatorname{tr}(\mathbf{S} \, \mathrm{d}\mathbf{U}) &= \operatorname{tr}\big(\mathbf{S}(\mathrm{d}\mathbf{X}\mathbf{1}_N)\mathbf{r}^T + \mathbf{S}\mathbf{r}(\mathrm{d}\mathbf{X}\mathbf{1}_N)^T\big) \notag\\
&= \operatorname{tr}\big(\mathbf{r}^T\mathbf{S}(\mathrm{d}\mathbf{X}\mathbf{1}_N)\big) + \operatorname{tr}\big((\mathrm{d}\mathbf{X}\mathbf{1}_N)^T\mathbf{S}\mathbf{r}\big) \notag\\
&= 2\operatorname{tr}\big(\mathbf{r}^T\mathbf{S}\,\mathrm{d}\mathbf{X}\mathbf{1}_N\big) \notag\\
&= 2\operatorname{tr}\big(\mathbf{1}_N\mathbf{r}^T\mathbf{S}\,\mathrm{d}\mathbf{X}\big) \notag\\
&= 2\operatorname{tr}\big((\mathbf{S}\mathbf{r}\mathbf{1}_N^T)^T \mathrm{d}\mathbf{X}\big).
\end{align}
Multiplying by the factor of $2$ from \eqref{eq:df_trace} yields the first gradient term:
\begin{equation}
\nabla_{\mathbf{X}}^{(1)} f = 4\mathbf{S}\mathbf{r}\mathbf{1}_N^T = 4\big(\mathbf{R} \circ \hat{\mathbf{A}}_G\big)\mathbf{r}\mathbf{1}_N^T.
\end{equation}

\textit{Part 2: Bilinear differential $\mathrm{d}(\mathbf{X}\hat{\mathbf{A}}_H\mathbf{X}^T)$.} By the product rule of matrix differentiation:
\begin{equation}
\mathrm{d}(\mathbf{X}\hat{\mathbf{A}}_H\mathbf{X}^T) = (\mathrm{d}\mathbf{X})\hat{\mathbf{A}}_H\mathbf{X}^T + \mathbf{X}\hat{\mathbf{A}}_H(\mathrm{d}\mathbf{X})^T.
\end{equation}
Substituting into \eqref{eq:df_trace}:
\begin{align}
\operatorname{tr}\big(\mathbf{R} \, \mathrm{d}(\mathbf{X}\hat{\mathbf{A}}_H\mathbf{X}^T)\big) &= \operatorname{tr}\big(\mathbf{R}(\mathrm{d}\mathbf{X})\hat{\mathbf{A}}_H\mathbf{X}^T\big) \notag\\
&= \operatorname{tr}\big(\hat{\mathbf{A}}_H\mathbf{X}^T\mathbf{R}\,\mathrm{d}\mathbf{X}\big) \notag\\
&\quad + \operatorname{tr}\big((\mathbf{R}\mathbf{X}\hat{\mathbf{A}}_H)^T\mathrm{d}\mathbf{X}\big) \notag\\
&= 2\operatorname{tr}\big((\mathbf{R}\mathbf{X}\hat{\mathbf{A}}_H)^T \mathrm{d}\mathbf{X}\big),
\end{align}
where we used the symmetry of $\mathbf{R}$ and $\hat{\mathbf{A}}_H$. Accounting for the minus sign in \eqref{eq:dR} and the factor of $2$ in \eqref{eq:df_trace}, the second gradient term is:
\begin{equation}
\nabla_{\mathbf{X}}^{(2)} f = -4\mathbf{R}\mathbf{X}\hat{\mathbf{A}}_H.
\end{equation}
Combining both parts:
\begin{equation}
\nabla_{\mathbf{X}} f(\mathbf{X}) = 4\big(\mathbf{R} \circ \hat{\mathbf{A}}_G\big)\mathbf{r}\mathbf{1}_N^T - 4\mathbf{R}\mathbf{X}\hat{\mathbf{A}}_H.
\end{equation}
To verify the rank-one broadcast property, note that $\mathbf{v} = (\mathbf{R} \circ \hat{\mathbf{A}}_G)\mathbf{r} \in \mathbb{R}^M$ is a column vector. The outer product $\mathbf{v}\mathbf{1}_N^T \in \mathbb{R}^{M \times N}$ replicates $v_i$ along all columns $j \in \{1, \ldots, N\}$. This completes the proof.
\end{proof}

\begin{corollary}[Evaluation Complexity of $\nabla_{\mathbf{X}} f$]
\label{cor:gradient_complexity}
The analytical gradient $\nabla_{\mathbf{X}} f(\mathbf{X})$ evaluates in strictly $\Theta(M^2N + MN^2) = \Theta(N^3)$ operations.
\end{corollary}

\begin{proof}
Let us enumerate the operations:
\begin{enumerate}
    \item Compute $\mathbf{r} = \mathbf{X}\mathbf{1}_N$: $\mathcal{O}(MN)$ operations.
    \item Compute $\mathbf{U} = \mathbf{r}\mathbf{r}^T$: $\mathcal{O}(M^2)$ operations.
    \item Compute $\mathbf{P} = \mathbf{X}\hat{\mathbf{A}}_H$: $\mathcal{O}(MN^2)$ operations.
    \item Compute $\mathbf{P}\mathbf{X}^T = \mathbf{X}\hat{\mathbf{A}}_H\mathbf{X}^T$: $\mathcal{O}(M^2N)$ operations.
    \item Compute $\mathbf{R} = (\mathbf{U} \circ \hat{\mathbf{A}}_G) - \mathbf{P}\mathbf{X}^T$: $\mathcal{O}(M^2)$ operations.
    \item Compute $\mathbf{v} = (\mathbf{R} \circ \hat{\mathbf{A}}_G)\mathbf{r}$: $\mathcal{O}(M^2)$ operations.
    \item Broadcast $\mathbf{v}\mathbf{1}_N^T$: $\mathcal{O}(MN)$ operations.
    \item Compute $\mathbf{R}\mathbf{P} = \mathbf{R}\mathbf{X}\hat{\mathbf{A}}_H$: $\mathcal{O}(M^2N)$ operations.
\end{enumerate}
All operations are bounded by $\mathcal{O}(M^2N + MN^2) = \Theta(N^3)$, which is strictly cubic and independent of finite-difference sampling.
\end{proof}

\subsection{Optional Shared-Scale Adjacency Normalization}
The solver requires symmetric adjacency matrices and defaults to using their input scales unchanged. It also supports two opt-in normalization modes. In \texttt{shared\_max} mode, both matrices are divided by the same scalar
\begin{equation}
s_{\max} = \max\big(\|\mathbf{A}_G\|_{\max},\|\mathbf{A}_H\|_{\max}\big), \qquad
\hat{\mathbf{A}}_G=\mathbf{A}_G/s_{\max},\quad \hat{\mathbf{A}}_H=\mathbf{A}_H/s_{\max}.
\label{eq:shared_scale}
\end{equation}
In \texttt{frobenius} mode, both matrices use the common geometric-mean scale
\begin{equation}
s_F = \sqrt{\|\mathbf{A}_G\|_F\,\|\mathbf{A}_H\|_F}, \qquad
\hat{\mathbf{A}}_G=\mathbf{A}_G/s_F,\quad \hat{\mathbf{A}}_H=\mathbf{A}_H/s_F.
\end{equation}
The implementation applies either scale only when the corresponding mode is selected and the scale is nonzero. A common scale preserves relative edge weights between the two inputs; choosing a normalization mode changes the objective scale and should therefore be reported with any experiment. The default is \texttt{none}.

\subsection{Frank-Wolfe Gap and Inner-Loop Stopping}

Within the GNCCP framework, the continuation path parameter $\zeta$ deforms the objective function according to:
\begin{equation}
J_\zeta(\mathbf{X}) = \begin{cases}
(1 - \zeta)F(\mathbf{X}) + \zeta\|\mathbf{X}\|_F^2, & \zeta \in [0, 1], \\
(1 + \zeta)F(\mathbf{X}) + \zeta\|\mathbf{X}\|_F^2, & \zeta \in [-1, 0).
\end{cases}
\label{eq:Jzeta_def}
\end{equation}
At each fixed $\zeta$, we solve $\min_{\mathbf{X} \in \mathcal{D}_L} J_\zeta(\mathbf{X})$ using the Frank-Wolfe conditional gradient method \cite{frank1956algorithm, jaggi2013revisiting}. At iterate $\mathbf{X}_t$, the Linear Minimization Oracle (LMO) solves:
\begin{equation}
\mathbf{Y}^* = \arg\min_{\mathbf{Y} \in \mathcal{D}_L} \big\langle \nabla J_\zeta(\mathbf{X}_t), \mathbf{Y} \big\rangle.
\label{eq:lmo_def}
\end{equation}

\begin{theorem}[Frank-Wolfe Gap Properties]
\label{thm:duality_gap}
Let the Frank-Wolfe duality gap at iterate $\mathbf{X}_t \in \mathcal{D}_L$ be defined as:
\begin{equation}
g_t = \big\langle \nabla J_\zeta(\mathbf{X}_t), \mathbf{X}_t - \mathbf{Y}^* \big\rangle.
\label{eq:gt_def}
\end{equation}
Then:
\begin{enumerate}
    \item $g_t \geq 0$ for all $\mathbf{X}_t \in \mathcal{D}_L$.
    \item If $J_\zeta$ is convex on $\mathcal{D}_L$, $g_t$ provides an upper bound on global suboptimality:
    \begin{equation}
    J_\zeta(\mathbf{X}_t) - J_\zeta(\mathbf{X}^*) \leq g_t.
    \end{equation}
    \item $g_t = 0$ if and only if $\mathbf{X}_t$ satisfies the first-order Karush-Kuhn-Tucker (KKT) stationary condition on $\mathcal{D}_L$.
\end{enumerate}
\end{theorem}

\begin{proof}
\textit{Proof of (1):} Since $\mathbf{Y}^*$ minimizes the linear form $\langle\nabla J_\zeta(\mathbf{X}_t), \mathbf{Y}\rangle$ over $\mathcal{D}_L$, and $\mathbf{X}_t$ is a feasible point in $\mathcal{D}_L$:
\begin{equation}
\langle\nabla J_\zeta(\mathbf{X}_t), \mathbf{Y}^*\rangle \leq \langle\nabla J_\zeta(\mathbf{X}_t), \mathbf{X}_t\rangle.
\end{equation}
Rearranging yields $g_t = \langle\nabla J_\zeta(\mathbf{X}_t), \mathbf{X}_t - \mathbf{Y}^*\rangle \geq 0$.

\textit{Proof of (2):} When $J_\zeta$ is convex, the first-order Taylor approximation forms a global underestimator everywhere on $\mathcal{D}_L$:
\begin{equation}
J_\zeta(\mathbf{X}^*) \geq J_\zeta(\mathbf{X}_t) + \big\langle \nabla J_\zeta(\mathbf{X}_t), \mathbf{X}^* - \mathbf{X}_t \big\rangle.
\end{equation}
Rearranging:
\begin{equation}
J_\zeta(\mathbf{X}_t) - J_\zeta(\mathbf{X}^*) \leq \big\langle \nabla J_\zeta(\mathbf{X}_t), \mathbf{X}_t - \mathbf{X}^* \big\rangle.
\end{equation}
Since $\mathbf{X}^* \in \mathcal{D}_L$, the definition of $\mathbf{Y}^*$ implies $\langle\nabla J_\zeta(\mathbf{X}_t), \mathbf{X}^*\rangle \geq \langle\nabla J_\zeta(\mathbf{X}_t), \mathbf{Y}^*\rangle$. Thus:
\begin{equation}
\big\langle \nabla J_\zeta(\mathbf{X}_t), \mathbf{X}_t - \mathbf{X}^* \big\rangle \leq \big\langle \nabla J_\zeta(\mathbf{X}_t), \mathbf{X}_t - \mathbf{Y}^* \big\rangle = g_t.
\end{equation}
Therefore, $J_\zeta(\mathbf{X}_t) - J_\zeta(\mathbf{X}^*) \leq g_t$.

\textit{Proof of (3):} The first-order KKT condition for constrained minimization over a convex set $\mathcal{D}_L$ requires:
\begin{equation}
\big\langle \nabla J_\zeta(\mathbf{X}_t), \mathbf{Y} - \mathbf{X}_t \big\rangle \geq 0, \quad \forall \mathbf{Y} \in \mathcal{D}_L.
\end{equation}
Setting $\mathbf{Y} = \mathbf{Y}^*$ yields $-\langle\nabla J_\zeta(\mathbf{X}_t), \mathbf{X}_t - \mathbf{Y}^*\rangle \geq 0 \implies -g_t \geq 0$. Since $g_t \geq 0$ from (1), we must have $g_t = 0$. Conversely, if $g_t = 0$, then for any $\mathbf{Y} \in \mathcal{D}_L$:
\begin{equation}
\langle\nabla J_\zeta(\mathbf{X}_t), \mathbf{Y} - \mathbf{X}_t\rangle \geq \langle\nabla J_\zeta(\mathbf{X}_t), \mathbf{Y}^* - \mathbf{X}_t\rangle = -g_t = 0,
\end{equation}
satisfying first-order stationarity. This completes the proof.
\end{proof}

\begin{remark}[Implementation of the Frank-Wolfe Gap]
At each inner iteration, v3 computes $g_t=\langle\nabla J_\zeta(\mathbf{X}_t),\mathbf{X}_t-\mathbf{Y}^*\rangle$ and exits that inner loop when $|g_t|<\texttt{tol}$. This check is used at every continuation level. It certifies an $\texttt{tol}$ global-suboptimality bound only when convexity of the current $J_\zeta$ has been established; at non-convex levels, a zero gap characterizes first-order stationarity, while a small nonzero gap is a stopping criterion rather than a global-optimality certificate. The loop is also bounded by \texttt{fw\_max\_iter}.
\end{remark}

\begin{remark}[Selection of Subgraph Cardinality $L$ in Virtual Screening]
In pairwise virtual screening and molecular property prediction, the common subgraph size $L$ between query graph $G$ and candidate library graph $H$ is adaptively assigned as $L = \max(1, \min(|V_G|, |V_H|) - 2)$. This allows up to two unassigned peripheral substituent atoms or scaffold variations while rigorously enforcing bijective matching over the conserved molecular core.
\end{remark}

\subsection{Relabeling Invariance Under Graph Isomorphisms}

A rigorous molecular kernel must be invariant to the arbitrary ordering of atom indices.

\begin{proposition}[Relabeling Invariance]
\label{prop:relabeling}
Let $\mathbf{P} \in \mathcal{P}_M$ and $\mathbf{Q} \in \mathcal{P}_N$ be full permutation matrices representing arbitrary relabelings of the atoms in graphs $G$ and $H$, such that $\mathbf{A}'_G = \mathbf{P}\mathbf{A}_G\mathbf{P}^T$, $\mathbf{A}'_H = \mathbf{Q}\mathbf{A}_H\mathbf{Q}^T$, and $\mathbf{C}' = \mathbf{P}\mathbf{C}\mathbf{Q}^T$. Then for any assignment matrix $\mathbf{X} \in \mathcal{D}_L$:
\begin{equation}
F\big(\mathbf{P}\mathbf{X}\mathbf{Q}^T; \mathbf{A}'_G, \mathbf{A}'_H, \mathbf{C}'\big) = F\big(\mathbf{X}; \mathbf{A}_G, \mathbf{A}_H, \mathbf{C}\big).
\label{eq:relabeling_invariance}
\end{equation}
\end{proposition}

\begin{proof}
Let $\mathbf{X}' = \mathbf{P}\mathbf{X}\mathbf{Q}^T$. For the linear trace term:
\begin{align}
\operatorname{tr}\big((\mathbf{C}')^T\mathbf{X}'\big) &= \operatorname{tr}\big(\mathbf{Q}\mathbf{C}^T\mathbf{P}^T\mathbf{P}\mathbf{X}\mathbf{Q}^T\big) \notag\\
&= \operatorname{tr}\big(\mathbf{Q}\mathbf{C}^T\mathbf{X}\mathbf{Q}^T\big) = \operatorname{tr}\big(\mathbf{C}^T\mathbf{X}\big),
\end{align}
by the cyclic property of trace and the fact that $\mathbf{P}^T\mathbf{P} = \mathbf{I}_M$ and $\mathbf{Q}^T\mathbf{Q} = \mathbf{I}_N$. For the bilinear term:
\begin{align}
\mathbf{X}'\hat{\mathbf{A}}'_H(\mathbf{X}')^T &= (\mathbf{P}\mathbf{X}\mathbf{Q}^T)(\mathbf{Q}\hat{\mathbf{A}}_H\mathbf{Q}^T)(\mathbf{Q}\mathbf{X}^T\mathbf{P}^T) \notag\\
&= \mathbf{P}(\mathbf{X}\hat{\mathbf{A}}_H\mathbf{X}^T)\mathbf{P}^T.
\end{align}
For the Hadamard term, using $\mathbf{r}' = \mathbf{X}'\mathbf{1}_N = \mathbf{P}\mathbf{X}\mathbf{Q}^T\mathbf{1}_N = \mathbf{P}\mathbf{X}\mathbf{1}_N = \mathbf{P}\mathbf{r}$, we have $\mathbf{U}' = \mathbf{r}'(\mathbf{r}')^T = \mathbf{P}\mathbf{r}\mathbf{r}^T\mathbf{P}^T = \mathbf{P}\mathbf{U}\mathbf{P}^T$. Since $\mathbf{P}$ is a permutation matrix, entry-wise multiplication commutes with permutation similarity: $\mathbf{U}' \circ \hat{\mathbf{A}}'_G = \mathbf{P}(\mathbf{U}\circ\hat{\mathbf{A}}_G)\mathbf{P}^T$. Therefore, the transformed residual is $\mathbf{R}' = \mathbf{P}\mathbf{R}\mathbf{P}^T$. Since the Frobenius norm is invariant under unitary and orthogonal transformations:
\begin{equation}
\|\mathbf{R}'\|_F^2 = \|\mathbf{P}\mathbf{R}\mathbf{P}^T\|_F^2 = \|\mathbf{R}\|_F^2.
\end{equation}
Thus, $F(\mathbf{X}')$ is strictly identical to $F(\mathbf{X})$.
\end{proof}

\subsection{Exact Closed-Form Quartic Line Search and Symmetric Matrix Identity}

In standard Frank-Wolfe optimization, each iteration updates $\mathbf{X}_{t+1} = \mathbf{X}_t - \lambda^* \mathbf{D}_{\text{FW}}$, where $\mathbf{D}_{\text{FW}} = \mathbf{X}_t - \mathbf{Y}^*$ is the descent direction and $\lambda^* \in [0, 1]$ is the step size. Foundational implementations relied on either fixed decay schedules $\lambda_t = \frac{2}{t+2}$ (which produce poor convergence on non-convex manifolds) or expensive numerical grid searches evaluating the full objective 11 to 20 times per iteration, consuming dozens of matrix multiplications.

Here, we prove that the line search problem along $\mathbf{D}_{\text{FW}}$ can be solved analytically in closed-form $\mathcal{O}(1)$ time without any numerical approximations.

\begin{theorem}[Closed-Form Quartic Line Search and Symmetric Matrix Identity]
\label{thm:quartic_line_search}
Let $\mathbf{X} \in \mathcal{D}_L$ and $\mathbf{D} = \mathbf{X} - \mathbf{Y}^* \in \mathbb{R}^{M \times N}$. Along the line $\mathbf{X}(\lambda) = \mathbf{X} - \lambda \mathbf{D}$ for $\lambda \in [0, 1]$, the homotopy objective $J_\zeta(\mathbf{X}(\lambda))$ is an exact fourth-degree polynomial:
\begin{equation}
J_\zeta(\lambda) = c_4 \lambda^4 + c_3 \lambda^3 + c_2 \lambda^2 + c_1 \lambda + c_0.
\label{eq:quartic_poly}
\end{equation}
The structural residual decomposes into quadratic matrix components:
\begin{equation}
\mathbf{R}(\lambda) = \mathbf{R}_0 + \lambda \mathbf{R}_1 + \lambda^2 \mathbf{R}_2,
\end{equation}
where:
\begin{align}
\mathbf{R}_0 &= (\mathbf{r}_X \mathbf{r}_X^T) \circ \hat{\mathbf{A}}_G - \mathbf{X}\hat{\mathbf{A}}_H\mathbf{X}^T, \\
\mathbf{R}_1 &= -\big(\mathbf{r}_X \mathbf{d}_r^T + \mathbf{d}_r \mathbf{r}_X^T\big) \circ \hat{\mathbf{A}}_G + \big(\mathbf{K}^T + \mathbf{K}\big), \\
\mathbf{R}_2 &= (\mathbf{d}_r \mathbf{d}_r^T) \circ \hat{\mathbf{A}}_G - \mathbf{D}\hat{\mathbf{A}}_H\mathbf{D}^T,
\end{align}
with $\mathbf{r}_X = \mathbf{X}\mathbf{1}_N$, $\mathbf{d}_r = \mathbf{D}\mathbf{1}_N$, and $\mathbf{K} = (\mathbf{D}\hat{\mathbf{A}}_H)\mathbf{X}^T \in \mathbb{R}^{M \times M}$. Crucially, because $\hat{\mathbf{A}}_H$ is symmetric ($\hat{\mathbf{A}}_H = \hat{\mathbf{A}}_H^T$), the cross-bilinear term satisfies:
\begin{equation}
\mathbf{X}\hat{\mathbf{A}}_H\mathbf{D}^T = \big(\mathbf{D}\hat{\mathbf{A}}_H\mathbf{X}^T\big)^T = \mathbf{K}^T,
\label{eq:symm_gemm}
\end{equation}
eliminating one full matrix multiplication. The Frobenius norm $\|\mathbf{R}(\lambda)\|_F^2 = \sum_{k=0}^4 s_k \lambda^k$ expands with scalar coefficients:
\begin{align}
s_0 &= \langle \mathbf{R}_0, \mathbf{R}_0 \rangle, \quad s_1 = 2\langle \mathbf{R}_0, \mathbf{R}_1 \rangle, \notag\\
s_2 &= \langle \mathbf{R}_1, \mathbf{R}_1 \rangle + 2\langle \mathbf{R}_0, \mathbf{R}_2 \rangle, \notag\\
s_3 &= 2\langle \mathbf{R}_1, \mathbf{R}_2 \rangle, \quad s_4 = \langle \mathbf{R}_2, \mathbf{R}_2 \rangle.
\end{align}
Letting $u_0 = \operatorname{tr}(\mathbf{C}^T\mathbf{X})$, $u_1 = -\operatorname{tr}(\mathbf{C}^T\mathbf{D})$, $x_0 = \|\mathbf{X}\|_F^2$, $x_1 = -2\langle \mathbf{X}, \mathbf{D} \rangle$, $x_2 = \|\mathbf{D}\|_F^2$, and $w = (1 - \zeta)$ for $\zeta \ge 0$ (or $1 + \zeta$ for $\zeta < 0$), the polynomial coefficients of $J_\zeta(\lambda)$ are:
\begin{align}
c_0 &= w(\alpha s_0 + (1-\alpha)u_0) + \zeta x_0, \notag\\
c_1 &= w(\alpha s_1 + (1-\alpha)u_1) + \zeta x_1, \notag\\
c_2 &= w\alpha s_2 + \zeta x_2, \quad c_3 = w\alpha s_3, \quad c_4 = w\alpha s_4.
\end{align}
Once the polynomial coefficients have been computed, the global minimizer $\lambda^* = \arg\min_{\lambda \in [0, 1]} J_\zeta(\lambda)$ is determined by a constant-size scalar calculation: evaluate $J_\zeta(\lambda)$ at $\{0, 1\}$ and at the real roots in $[0, 1]$ of the cubic derivative:
\begin{equation}
J'_\zeta(\lambda) = 4c_4 \lambda^3 + 3c_3 \lambda^2 + 2c_2 \lambda + c_1 = 0,
\label{eq:cubic_root}
\end{equation}
which are solved analytically via Cardano's formula.
\end{theorem}

\begin{proof}
Substituting $\mathbf{X}(\lambda) = \mathbf{X} - \lambda \mathbf{D}$ into $\mathbf{r}(\lambda) = \mathbf{X}(\lambda)\mathbf{1}_N = \mathbf{r}_X - \lambda \mathbf{d}_r$, the co-matching indicator matrix becomes:
\begin{align}
\mathbf{U}(\lambda) &= \mathbf{r}(\lambda)\mathbf{r}(\lambda)^T \notag\\
&= (\mathbf{r}_X - \lambda \mathbf{d}_r)(\mathbf{r}_X - \lambda \mathbf{d}_r)^T \notag\\
&= \mathbf{r}_X \mathbf{r}_X^T - \lambda(\mathbf{r}_X \mathbf{d}_r^T + \mathbf{d}_r \mathbf{r}_X^T) + \lambda^2 \mathbf{d}_r \mathbf{d}_r^T.
\end{align}
Similarly, the projected adjacency matrix expands as:
\begin{align}
\mathbf{X}(\lambda)\hat{\mathbf{A}}_H\mathbf{X}(\lambda)^T &= (\mathbf{X} - \lambda \mathbf{D})\hat{\mathbf{A}}_H(\mathbf{X} - \lambda \mathbf{D})^T \notag\\
&= \mathbf{X}\hat{\mathbf{A}}_H\mathbf{X}^T - \lambda\big(\mathbf{X}\hat{\mathbf{A}}_H\mathbf{D}^T + \mathbf{D}\hat{\mathbf{A}}_H\mathbf{X}^T\big) \notag\\
&\quad + \lambda^2 \mathbf{D}\hat{\mathbf{A}}_H\mathbf{D}^T.
\end{align}
Since $\hat{\mathbf{A}}_H = \hat{\mathbf{A}}_H^T$, we have $(\mathbf{X}\hat{\mathbf{A}}_H\mathbf{D}^T)^T = \mathbf{D}\hat{\mathbf{A}}_H^T\mathbf{X}^T = \mathbf{D}\hat{\mathbf{A}}_H\mathbf{X}^T = \mathbf{K}$. Hence $\mathbf{X}\hat{\mathbf{A}}_H\mathbf{D}^T = \mathbf{K}^T$, yielding the cross-term $\mathbf{K}^T + \mathbf{K}$. Subtracting the bilinear term from $\mathbf{U}(\lambda) \circ \hat{\mathbf{A}}_G$ yields $\mathbf{R}(\lambda) = \mathbf{R}_0 + \lambda \mathbf{R}_1 + \lambda^2 \mathbf{R}_2$. Expanding the trace inner product:
\begin{align}
\|\mathbf{R}(\lambda)\|_F^2 &= \operatorname{tr}\big((\mathbf{R}_0 + \lambda \mathbf{R}_1 + \lambda^2 \mathbf{R}_2)^T \notag\\
&\qquad \cdot (\mathbf{R}_0 + \lambda \mathbf{R}_1 + \lambda^2 \mathbf{R}_2)\big) \notag\\
&= \|\mathbf{R}_0\|_F^2 + 2\lambda\langle\mathbf{R}_0, \mathbf{R}_1\rangle \notag\\
&\quad + \lambda^2\big(\|\mathbf{R}_1\|_F^2 + 2\langle\mathbf{R}_0, \mathbf{R}_2\rangle\big) \notag\\
&\quad + 2\lambda^3\langle\mathbf{R}_1, \mathbf{R}_2\rangle + \lambda^4\|\mathbf{R}_2\|_F^2.
\end{align}
Linearity of the trace inner product ensures the unary term is affine in $\lambda$, while the quadratic Frobenius norm $\|\mathbf{X}(\lambda)\|_F^2$ expands to $x_0 + x_1 \lambda + x_2 \lambda^2$. Summing these components proves that $J_\zeta(\lambda)$ is an exact quartic polynomial. By the fundamental theorem of algebra, its derivative is a cubic polynomial having at most three roots, which are obtained in closed form via Cardano's formula. Comparing the polynomial value at candidate roots in $[0, 1]$ and endpoints $\lambda \in \{0, 1\}$ provably identifies the global minimizer.
\end{proof}

\begin{remark}[Computational Impact of Theorem~\ref{thm:quartic_line_search}]
Grid-search implementations evaluate the objective at multiple candidate step sizes, potentially repeating matrix products. The baseline implementation in this project uses an 11-point grid. V3 instead computes the quartic coefficients once per Frank-Wolfe step, reuses products cached during the gradient calculation, and evaluates the resulting scalar polynomial at its endpoints and real stationary points. The polynomial minimum has no grid discretization error; computing its coefficients still requires matrix operations.
\end{remark}

% =========================================================
% V. ALGORITHMIC ARCHITECTURE AND REALIZATION
% =========================================================

\section{Algorithmic Architecture and Realization}
\label{sec:algorithms}

The Fast-WCS architecture couples analytical homotopy continuation with exact linear assignment and 3D geometric superposition.

Unless otherwise specified, v3 uses $\alpha=0.7$, cosine continuation with adaptive level changes disabled, $8$ continuation levels, at most $3$ Frank-Wolfe steps per level, tolerance $10^{-3}$, float32 solver matrices, quartic line search, the exact augmented oracle, one uniform-barycenter start, and local refinement enabled with a maximum of $20$ rounds. Multiple restarts are optional; after the uniform start, additional starts are seeded random feasible assignment vertices, and the run with the lowest relaxed objective is retained before discretization.

\begin{algorithm}[t]
\caption{Fast-WCS Analytical Homotopy and Frank-Wolfe Optimization (Phase 1)}
\label{alg:fast_wcs}
\begin{algorithmic}[1]
\REQUIRE Symmetric adjacency matrices $\mathbf{A}_G \in \mathbb{R}^{M \times M}$ and $\mathbf{A}_H \in \mathbb{R}^{N \times N}$; optional unary cost $\mathbf{C} \in \mathbb{R}^{M \times N}$; target size $L$; parameters $\alpha$, tolerance, continuation levels $K$, and inner budget $T_{\text{FW}}$.
\ENSURE Binary correspondence $\mathbf{X}^* \in \mathcal{P}_L$, matched pairs $\mathcal{M}$, and objective and similarity scores.
\STATE Validate adjacency shapes and symmetry; set $\mathbf{C}=\mathbf{0}$ if no unary cost is supplied
\STATE Optionally apply the configured common adjacency scale (default: no normalization)
\STATE Initialize $\mathbf{X} \leftarrow \frac{L}{MN}\mathbf{1}_{M \times N}$
\FOR{$k = 0$ to $K - 1$}
    \STATE $\zeta \leftarrow \cos\big(\frac{\pi k}{K - 1}\big)$ \COMMENT{Cosine Continuation Schedule}
    \FOR{$t = 1$ to $T_{\text{FW}}$}
        \STATE Compute $\mathbf{r}=\mathbf{X}\mathbf{1}_N$, $\mathbf{P}=\mathbf{X}\mathbf{A}_H$, and $\mathbf{R}=(\mathbf{r}\mathbf{r}^T\circ\mathbf{A}_G)-\mathbf{P}\mathbf{X}^T$
        \STATE $\nabla_{\mathbf{X}} f \leftarrow 4[(\mathbf{R}\circ\mathbf{A}_G)\mathbf{r}]\mathbf{1}_N^T - 4\mathbf{R}\mathbf{P}$ \COMMENT{Theorem~\ref{thm:gradient}}
        \STATE $\nabla F \leftarrow \alpha \nabla_{\mathbf{X}} f + (1 - \alpha)\mathbf{C}$
        \STATE $\nabla J_\zeta \leftarrow (1 \mp \zeta)\nabla F + 2\zeta \mathbf{X}$ \COMMENT{Continuation gradient}
        \STATE $\mathbf{Y}^* \leftarrow \text{AugmentedHungarianLMO}(\nabla J_\zeta, L, M, N)$
        \STATE $g_t \leftarrow \langle \nabla J_\zeta, \mathbf{X} - \mathbf{Y}^* \rangle$
        \IF{$|g_t| < \texttt{tol}$}
            \STATE \textbf{break} \COMMENT{Inner-loop stopping test; not a global certificate at non-convex levels}
        \ENDIF
        \STATE $\mathbf{D}_{\text{FW}} \leftarrow \mathbf{X} - \mathbf{Y}^*$ \COMMENT{Descent direction}
        \STATE $\lambda^* \leftarrow \text{ExactQuarticSearch}(\mathbf{X}, \mathbf{Y}^*, \text{cached products}, \mathbf{A}_G, \mathbf{A}_H, \mathbf{C}, \zeta)$
        \STATE $\mathbf{X} \leftarrow (1 - \lambda^*)\mathbf{X} + \lambda^* \mathbf{Y}^*$ \COMMENT{Convex step in $\mathcal{D}_L$}
    \ENDFOR
\ENDFOR
\STATE $\mathbf{X}^* \leftarrow \text{AugmentedHungarianLMO}(-\mathbf{X}, L, M, N)$
\STATE If enabled, refine $\mathbf{X}^*$ with vectorized target-replacement, exchange, and source-replacement moves up to the configured round limit
\STATE Extract matched atom pairs $\mathcal{M} = \{(i, j) \mid X^*_{ij} = 1\}$
\STATE $F^* \leftarrow \alpha \|\mathbf{U}^* \circ \hat{\mathbf{A}}_G - \mathbf{X}^*\hat{\mathbf{A}}_H(\mathbf{X}^*)^T\|_F^2 + (1 - \alpha)\operatorname{tr}(\mathbf{C}^T\mathbf{X}^*)$
\STATE $k(G,H) \leftarrow \exp(-\gamma F^* / \max(L,1))$
\RETURN $\mathbf{X}^*$, $\mathcal{M}$, $F^*$, $k(G,H)$
\end{algorithmic}
\end{algorithm}

\begin{algorithm}[t]
\caption{Phase 2 Closed-Form 3D Kabsch Superposition}
\label{alg:kabsch}
\begin{algorithmic}[1]
\REQUIRE Matched pairs $\mathcal{M} = \{(i_l, j_l)\}_{l=1}^L$; coordinate matrices $\mathbf{P}_{\text{all}} \in \mathbb{R}^{M \times 3}$ and $\mathbf{Q}_{\text{all}} \in \mathbb{R}^{N \times 3}$.
\ENSURE Rigid rotation $\mathbf{R}^* \in \mathrm{SO}(3)$, translation $\mathbf{t}^* \in \mathbb{R}^3$, and RMSD over all matched pairs.
\STATE Extract matched matrices $\mathbf{P}_L, \mathbf{Q}_L \in \mathbb{R}^{L \times 3}$
\STATE Compute centroids $\boldsymbol{\mu}_P = \frac{1}{L}\mathbf{1}_L^T\mathbf{P}_L, \; \boldsymbol{\mu}_Q = \frac{1}{L}\mathbf{1}_L^T\mathbf{Q}_L$
\STATE Center coordinates $\bar{\mathbf{P}}_L = \mathbf{P}_L - \mathbf{1}_L\boldsymbol{\mu}_P, \; \bar{\mathbf{Q}}_L = \mathbf{Q}_L - \mathbf{1}_L\boldsymbol{\mu}_Q$
\STATE Compute $3 \times 3$ cross-covariance $\mathbf{H} = \bar{\mathbf{P}}_L^T \bar{\mathbf{Q}}_L \in \mathbb{R}^{3 \times 3}$ \COMMENT{$\mathcal{O}(L)$ work}
\STATE Compute SVD: $\mathbf{H} = \mathbf{W}\mathbf{\Sigma}\mathbf{V}^T$ \COMMENT{$\mathcal{O}(1)$ work}
\STATE $d \leftarrow \det(\mathbf{V}\mathbf{W}^T)$ \COMMENT{Reflection detector: $+1$ or $-1$}
\STATE $\mathbf{R}^* \leftarrow \mathbf{V} \operatorname{diag}(1, 1, d) \mathbf{W}^T \in \mathrm{SO}(3)$ \COMMENT{Optimal Kabsch rotation}
\STATE $\mathbf{t}^* \leftarrow \boldsymbol{\mu}_Q^T - \mathbf{R}^*\boldsymbol{\mu}_P^T$ \COMMENT{Optimal translation}
\STATE $\text{RMSD} \leftarrow \sqrt{\frac{1}{L}\sum_{l=1}^L \|\mathbf{R}^* \bar{\mathbf{P}}_{L, l}^T - \bar{\mathbf{Q}}_{L, l}^T\|_2^2}$
\RETURN $\mathbf{R}^*, \mathbf{t}^*, \text{RMSD}$
\end{algorithmic}
\end{algorithm}

\subsection{Cosine Continuation Annealing Schedule}

The v3 solver uses a fixed cosine continuation schedule by default. For $K>1$ levels, it is:
\begin{equation}
\zeta_k = \cos\!\left(\frac{\pi k}{K - 1}\right), \quad k \in \{0, 1, \ldots, K - 1\}.
\label{eq:cosine_schedule}
\end{equation}
The schedule uses smaller changes near $\zeta=\pm1$ and larger changes near $\zeta=0$. A linear schedule is also available. Adaptive insertion or skipping of levels is an optional setting and is disabled by default; the default run visits the prescribed levels and applies the gap stopping test only within each level's Frank-Wolfe loop.

\subsection{Exact Burkard Augmented Hungarian Linear Minimization Oracle}
\label{subsec:burkard}

The Linear Minimization Oracle in \eqref{eq:lmo_def} requires solving a partial linear sum assignment problem over the continuous polytope $\mathcal{D}_L$:
\begin{equation}
\min_{\mathbf{Y} \in \mathcal{D}_L} \operatorname{tr}(\mathbf{M}^T \mathbf{Y}) = \min_{\mathbf{Y} \in \mathcal{D}_L} \sum_{i=1}^M \sum_{j=1}^N M_{ij} Y_{ij},
\label{eq:lmo_problem}
\end{equation}
where $\mathbf{M} = \nabla J_\zeta(\mathbf{X}_t) \in \mathbb{R}^{M \times N}$. When $L < \min(M, N)$, standard rectangular Hungarian matching cannot enforce the exact cardinality constraint $\sum_{i, j} Y_{ij} = L$. In prior work, Yang et al.~\cite{yang2018algorithm} settled for solving a full $M$-assignment and heuristically dropping the $M - L$ entries with smallest gradient (RLX1F). However, when $\nabla J_\zeta$ contains negative entries, this heuristic truncation misidentifies dummy nodes, frequently leading to poor descent directions.

To eliminate this lossy heuristic, Fast-WCS formulates an exact augmented Hungarian matrix based on the Burkard assignment construction \cite{loiola2007}. We construct an augmented square cost matrix $\tilde{\mathbf{M}} \in \mathbb{R}^{(M+N-L) \times (M+N-L)}$:
\begin{equation}
\tilde{\mathbf{M}} = \begin{bmatrix}
\mathbf{M}_{M \times N} & \mathbf{0}_{M \times (M - L)} \\
\mathbf{0}_{(N - L) \times N} & \infty \cdot \mathbf{1}_{(N - L) \times (M - L)}
\end{bmatrix}.
\label{eq:augmented_hungarian}
\end{equation}
Because the bottom-right slack block carries an infinite penalty ($\infty$), the optimal bijection on $\tilde{\mathbf{M}}$ cannot pair any dummy row with a dummy column. Thus, all $(N - L)$ dummy rows must match into the $N$ real columns (at zero cost), leaving exactly $N - (N - L) = L$ real columns to match with real rows. Symmetrically, all $(M - L)$ dummy columns must match into the $M$ real rows (at zero cost), leaving exactly $M - (M - L) = L$ real rows to match with real columns.

Crucially, because every feasible matching on $\tilde{\mathbf{M}}$ contains exactly $L$ real-real assignments and zero-cost slack assignments, any uniform constant shift of $\mathbf{M}$ shifts all feasible assignment costs equally by $L \cdot c$, leaving the optimal assignment completely invariant. Thus, no artificial positivity shifts or finite big-$M$ heuristics are required. Extracting the entries where both row $< M$ and col $< N$ produces an assignment matrix $\mathbf{Y}^* \in \mathcal{P}_L$ with strictly $\sum_{i, j} Y^*_{ij} = L$, solved in strictly $\mathcal{O}((M + N - L)^3) = \Theta(N^3)$ operations by standard Kuhn-Munkres algorithms with zero heuristic truncation error.

\subsection{Vectorized BLAS-3 Combinatorial 2-Opt Local Refinement}
\label{subsec:2opt}

While the continuous GNCCP homotopy reliably reaches high-quality attraction basins, continuous relaxation followed by binary projection can leave the discrete mapping sub-optimal with respect to local pairwise vertex swaps. Fast-WCS incorporates a discrete combinatorial 2-opt search comprising three move families:
\begin{enumerate}
    \item \textbf{Target Column Replacement (Move 1):} A matched row $i$ replaces its currently assigned column $c_i$ with an unmatched column $u \in \mathcal{U}_H$.
    \item \textbf{Pairwise Target Exchange (Move 2):} Two matched rows $i$ and $j$ swap their assigned target columns ($c_i \leftrightarrow c_j$).
    \item \textbf{Source Row Replacement (Move 3):} An unmatched source row $v \in \mathcal{U}_G$ replaces a matched row $i$.
\end{enumerate}

\begin{proposition}[Level-3 BLAS Vectorized 2-Opt Refinement]
\label{prop:2opt}
Let $\mathbf{P} = \hat{\mathbf{A}}_G[\mathrm{rows}, \mathrm{rows}] \in \mathbb{R}^{L \times L}$ and $\mathbf{Q} = \hat{\mathbf{A}}_H[\mathrm{cols}, \mathrm{cols}] \in \mathbb{R}^{L \times L}$ denote the submatrices corresponding to currently matched vertices. The objective deltas $\Delta F$ for all $L(L-1)/2$ pairwise target exchanges, all $L(N-L)$ target replacements, and all $L(M-L)$ source replacements can be evaluated simultaneously using Level-3 BLAS matrix multiplications ($\mathbf{P}\mathbf{Q}^T$, $\mathbf{P}\mathbf{B}^T$, and $\mathbf{B}\mathbf{Q}^T$) in $\mathcal{O}(L^2 N)$ arithmetic operations, eliminating nested element-wise loops.
\end{proposition}

\noindent\textbf{Vectorized Scoring via Expanded-Square Identity:} Fast-WCS eliminates all Python-level pair loops by proving that all candidate moves can be scored simultaneously using matrix multiplication. Let $\mathbf{P} = \hat{\mathbf{A}}_G[\text{rows}, \text{rows}] \in \mathbb{R}^{L \times L}$ and $\mathbf{Q} = \hat{\mathbf{A}}_H[\text{cols}, \text{cols}] \in \mathbb{R}^{L \times L}$ denote the submatrices corresponding to the currently matched vertices. Expanding the squared Frobenius difference:
\begin{equation}
\sum_{k=1}^L (P_{ak} - Q_{bk})^2 = \sum_{k=1}^L P_{ak}^2 + \sum_{k=1}^L Q_{bk}^2 - 2 \sum_{k=1}^L P_{ak} Q_{bk}.
\end{equation}
Crucially, the all-pairs cross term across all candidate combinations is computed simultaneously by the single Level-3 BLAS matrix product $\mathbf{P}\mathbf{Q}^T$:
\begin{equation}
\mathbf{E} = (\mathbf{P} \odot \mathbf{P})\mathbf{1}_L \mathbf{1}_L^T + \mathbf{1}_L \mathbf{1}_L^T (\mathbf{Q} \odot \mathbf{Q})^T - 2 \mathbf{P}\mathbf{Q}^T \in \mathbb{R}^{L \times L}.
\end{equation}
Let $\mathbf{p}_d = \operatorname{diag}(\mathbf{P})$ and $\mathbf{q}_d = \operatorname{diag}(\mathbf{Q})$. For Move 2 (pairwise exchange), the self-interaction diagonals ($k \in \{i, j\}$) are subtracted analytically via matrix broadcasts:
\begin{align}
\mathbf{S}_1 &= \mathbf{E} - (\mathbf{p}_d\mathbf{1}_L^T - \mathbf{Q}^T)^{\odot 2} - (\mathbf{P} - \mathbf{1}_L\mathbf{q}_d^T)^{\odot 2}, \\
\Delta_{\text{struct}} &= 2\big( (\mathbf{S}_1 - \mathbf{S}_{1b}) + (\mathbf{S}_2 - \mathbf{S}_{2b}) \big).
\end{align}
Similarly, for Move 1, let $\mathbf{B} = \hat{\mathbf{A}}_H[\text{unmatched}, \text{cols}] \in \mathbb{R}^{U \times L}$. All $L \times U$ candidate column replacements are scored simultaneously using $\mathbf{P}\mathbf{B}^T$. For Move 3, candidate source row replacements are scored using $\mathbf{B}\mathbf{Q}^T$.
For each family, the best candidate is selected by a vectorized minimum over its delta matrix. Since the refinement has a configurable round cap, termination at that cap does not by itself imply that no improving move remains.

\subsection{Optional 3D Spatial Hungarian Polish}
\label{subsec:dgm}

When physical 3D coordinates are present, the point-cloud adapter first builds symmetric Gaussian distance adjacencies and solves the graph-matching problem. If spatial polishing is enabled (the adapter's \texttt{n\_dgm\_iters} value is greater than one) and the initial all-pair RMSD is greater than $0.10\,\text{\AA}$, it performs one spatial polish:
\begin{enumerate}
    \item The matched coordinates are rigidly aligned by Kabsch, producing $\mathbf{R}^*$, $\mathbf{t}^*$, and the RMSD over all matched pairs.
    \item If the RMSD is above the threshold, coordinates from the first graph are transformed as $\mathbf{P}_{\text{trans}} = \mathbf{P}\mathbf{R}^{*T} + \mathbf{t}^*$. A pairwise spatial Euclidean cost matrix is computed and divided by its maximum value when nonzero:
    \begin{equation}
    C_{\text{spatial}, ij} = \frac{\|\mathbf{p}_{\text{trans}, i} - \mathbf{q}_j\|_2}{\max_{a, b} \|\mathbf{p}_{\text{trans}, a} - \mathbf{q}_b\|_2},
    \end{equation}
    If an appearance/feature cost $\mathbf{C}$ was supplied, it is blended equally with the spatial cost; otherwise the spatial cost alone is used. The augmented Hungarian assignment is applied directly to this cost, followed by local refinement when enabled. The candidate replaces the original result only if its all-pair Kabsch RMSD is strictly lower. This step does not run another GNCCP continuation pass.
\end{enumerate}

\subsection{Phase 2 Closed-Form 3D Kabsch Superposition}

The discrete assignment $\mathbf{X}^* \in \mathcal{P}_L$ specifies the matched pairs $\mathcal{M} = \{(i_l, j_l)\}_{l=1}^L$. The Kabsch routine fits a proper rigid rotation and translation to these pairs and reports RMSD over all $L$ pairs. In v3, trimming is disabled for this routine by default; any optional trimmed-fit diagnostic is separate from the reported all-pair RMSD.

As detailed in Algorithm~\ref{alg:kabsch}, let $\mathbf{P}_L$ and $\mathbf{Q}_L$ contain the matched coordinates. Centering and forming the $3 \times 3$ cross-covariance matrix $\mathbf{H}=\bar{\mathbf{P}}_L^T\bar{\mathbf{Q}}_L$ require $\mathcal{O}(L)$ work. The SVD is fixed-size in 3D, and the determinant correction enforces a proper rotation with determinant $+1$. The implementation returns the transform and RMSD evaluated over every matched pair.

\subsection{Comprehensive Complexity Comparison}

Table~\ref{tab:complexity_comparison} compares theoretical and empirical profiles of Fast-WCS and the listed graph-matching methods. For one Frank-Wolfe step, v3 evaluates the gradient in $\mathcal{O}(MN^2+M^2N)$ arithmetic and solves an augmented assignment of size $M+N-L$. When $M$ and $N$ are of the same order, both terms are cubic in graph size. Total runtime also depends on the number of continuation levels, inner steps, restarts, and local-refinement rounds; the non-convex WCS problem is not solved globally by these bounds.

\begin{table*}[t]
\centering
\caption{Asymptotic Computational and Memory Complexity Profiles}
\label{tab:complexity_comparison}
\fontsize{7.2}{8.6}\selectfont
\setlength{\tabcolsep}{3.0pt}
\begin{tabular*}{\textwidth}{@{\extracolsep{\fill}}lcccc@{}}
\toprule
\textbf{Method Framework} & \textbf{Space Complexity} & \textbf{Objective / Gradient Evaluation} & \textbf{LMO / Assignment Solver} & \textbf{Total Complexity Per Iteration} \\
\midrule
Spectral Matching (SM) \cite{leordeanu2005} & $\mathcal{O}(M^2 N^2)$ & Power Iteration & Greedy $\mathcal{O}(MN)$ & $\mathcal{O}(N^4)$ \\
Graduated Assignment (GA) \cite{gold1996} & $\mathcal{O}(M^2 N^2)$ & Matrix-Vector Product & Softmax Normalization $\mathcal{O}(MN)$ & $\mathcal{O}(N^4)$ \\
Projected Gradient (PGM) \cite{zaslavskiy2009} & $\mathcal{O}(M^2 N^2)$ & Affinity Step & Projection $\mathcal{O}(N^3)$ & $\mathcal{O}(N^4)$ \\
\midrule
GNCCP RLX1 \cite{yang2018algorithm} & $\mathcal{O}(N^2)$ & 4th-Order Polynomial Expansion & LP Solver $\mathcal{O}(N^6)$ & $\mathcal{O}(N^{5.4})$ (empirical) \\
GNCCP RLX2 \cite{yang2018algorithm} & $\mathcal{O}(N^2)$ & 6th-Order Polynomial Expansion & LP Solver $\mathcal{O}(N^6)$ & $\mathcal{O}(N^{5.4})$ (empirical) \\
GNCCP RLX1F \cite{yang2018algorithm} & $\mathcal{O}(N^2)$ & 4th-Order Polynomial Expansion & Truncation $\mathcal{O}(N^3)$ & $\mathcal{O}(N^{3.5})$ (lossy heuristic) \\
Naive Numerical Baseline 2018 & $\mathcal{O}(N^2)$ & Finite-Difference Sampling $\Theta(N^5)$ & Hungarian $\Theta(N^3)$ & $\Theta(N^5)$ \\
\midrule
\textbf{Fast-WCS v3} & $\mathcal{O}(M^2+N^2+MN)$ & \textbf{Exact analytical gradient $\mathcal{O}(MN^2+M^2N)$} & \textbf{Augmented Hungarian $\mathcal{O}((M+N-L)^3)$} & $\mathcal{O}(MN^2+M^2N+(M+N-L)^3)$ per FW step \\
\bottomrule
\end{tabular*}
\end{table*}

% =========================================================
% VI. REPLICATED SYNTHETIC BENCHMARK EVALUATION
% =========================================================

\begin{figure*}[t]
\centering
\includegraphics[width=0.98\textwidth,height=1.40in,keepaspectratio]{fig_wcs_synthetic_benchmark.pdf}
\caption{Comprehensive synthetic graph matching benchmark comparison replicating the protocol of Yang et al.~\cite{yang2018algorithm} across 50 Monte Carlo trials per setting. Top row (a--d): Matching accuracy comparison with respect to noise level $\sigma$, problem size $M$, outlier number $(M - L)$, and edge density $\rho$. Bottom row (e--h): Running time comparison and empirical complexity scaling slopes across the same four scenarios. Fast-WCS consistently achieves superior matching accuracy across all regimes while maintaining a strictly cubic empirical complexity slope of $\mathbf{\Theta(N^{3.0})}$, running up to two orders of magnitude faster than prior GNCCP relaxations ($\mathcal{O}(N^{5.4})$).}
\label{fig:synthetic_benchmark}
\end{figure*}

\section{Replicated Synthetic Benchmark Evaluation}
\label{sec:synthetic_experiments}

To validate the theoretical advantages of Fast-WCS against established state-of-the-art baselines, we replicated the benchmark protocol established by Yang et al. in Section~III-A of their IEEE TNNLS paper \cite{yang2018algorithm}.

\subsection{Experimental Settings and Protocol}

In this protocol, two spatial point sets $\mathcal{G} = \{\mathbf{g}_i\}_{i=1}^M$ and $\mathcal{H} = \{\mathbf{h}_j\}_{j=1}^N$ are uniformly sampled in the unit square $[0, 1]^2$. A ground-truth partial permutation matrix $\mathbf{X}_{\text{gt}} \in \mathcal{P}_L$ containing exactly $L$ ones is generated uniformly at random. A subset of points in $\mathcal{G}$ are created by permuting points in $\mathcal{H}$ with additive Gaussian spatial noise:
\begin{equation}
\mathbf{g}_i = \mathbf{h}_j + \boldsymbol{\eta}, \quad \boldsymbol{\eta} \sim \mathcal{N}(\mathbf{0}, \sigma^2 \mathbf{I}_2), \quad \text{if } X_{\text{gt}, ij} = 1.
\label{eq:spatial_noise}
\end{equation}
Pairwise Euclidean distances determine edge weights. Sparse adjacency matrices $\mathbf{A}_G$ and $\mathbf{A}_H$ are constructed by thresholding edges to achieve a target edge density $\rho \in [0.1, 1.0]$. Structural noise is introduced by randomly adding and deleting $(1/2)\sigma \cdot \#\text{Edge}$ edges in each graph.

We benchmark seven algorithms: (1) \textbf{Fast-WCS} (our proposed analytical framework); (2) \textbf{RLX1} (first polynomial relaxation of Yang et al. \cite{yang2018algorithm}); (3) \textbf{RLX1F} (RLX1 with fast heuristic Hungarian truncation \cite{yang2018algorithm}); (4) \textbf{RLX2} (second polynomial relaxation \cite{yang2018algorithm}); (5) \textbf{PGM} (Projected Gradient Matching \cite{zaslavskiy2009}); (6) \textbf{GA} (Graduated Assignment \cite{gold1996}); and (7) \textbf{SM} (Spectral Matching \cite{leordeanu2005}). Matching accuracy is quantified as:
\begin{equation}
\text{Accuracy} = \frac{\operatorname{tr}(\mathbf{X}_{\text{gt}}^T \mathbf{X}^*)}{L} \times 100\%.
\end{equation}
Each setting is evaluated across 50 independent Monte Carlo runs.

\subsection{Experimental Results Across Four Scenarios}

The matching accuracy comparisons across the four benchmark scenarios are plotted in Fig.~\ref{fig:synthetic_benchmark}(a--d):

\begin{itemize}
    \item \textbf{Scenario 1: Noise Level $\sigma$ (Fig.~\ref{fig:synthetic_benchmark}a):} Problem size is fixed at $M = 30$, $N = 35$, $L = 25$, edge density $\rho = 0.5$, while spatial noise $\sigma$ increases from $0.00$ to $0.10$. Fast-WCS achieves superior matching accuracy across all regimes ($89.3\%$ at $\sigma = 0.00$ in $5.4\,\text{ms}$, $52.0\%$ at $\sigma = 0.02$ in $3.4\,\text{ms}$, $44.0\%$ at $\sigma = 0.04$ in $3.2\,\text{ms}$, and $30.7\%$ at $\sigma = 0.08$ in $6.9\,\text{ms}$). In sharp contrast, RLX1 fluctuates between $4.0\%$ and $21.3\%$, RLX1F drops sharply, RLX2 degrades, PGM remains below $10\%$, GA achieves $6.7\%\text{--}13.3\%$, and SM collapses to $2.7\%$. Fast-WCS demonstrates marked robustness because its exact analytical gradient eliminates artificial local extrema.
    \item \textbf{Scenario 2: Problem Size $M$ Extended to $60$ (Fig.~\ref{fig:synthetic_benchmark}b):} With $\sigma = 0.03$, $N = M + 5$, $L = M - 5$, $\rho = 0.5$, problem size $M$ scales from $20$ to $60$ ($N=65, L=55$). Despite the factorial expansion of the combinatorial search space ($\binom{M}{L}\binom{N}{L}L!$), Fast-WCS consistently maintains the highest matching accuracy across all sizes ($53.3\%$ at $M=20$, $53.8\%$ at $M=36$, and $59.4\%$ at $M=60$), executing in strictly sub-$25\,\text{ms}$ wall-clock time ($6.3\,\text{ms}$ at $M=20$, $10.9\,\text{ms}$ at $M=44$, and $23.1\,\text{ms}$ at $M=60$). In contrast, RLX1 fluctuates between $6.8\%$ and $37.6\%$ while requiring $46.3\,\text{ms}$ at $M=60$, and RLX2 collapses.
    \item \textbf{Scenario 3: Outlier Count $(M - L)$ (Fig.~\ref{fig:synthetic_benchmark}c):} With $\sigma = 0.03$, $M = 30$, $N = 35$, $\rho = 0.5$, the number of outlier nodes $(M - L)$ increases from $0$ to $10$. When no outliers exist ($M = L = 30$), Fast-WCS achieves $60.0\%$ in $6.5\,\text{ms}$, reaching $76.2\%$ with $2$ outliers ($7.9\,\text{ms}$) and $69.2\%$ with $4$ outliers ($7.7\,\text{ms}$). Even under 10 outlier distractors, Fast-WCS maintains $46.7\%$ accuracy in $5.3\,\text{ms}$, whereas RLX1 collapses to $2.8\%\text{--}15.2\%$, RLX1F drops to $3.3\%$, and SM collapses to $5.0\%$.
    \item \textbf{Scenario 4: Edge Density $\rho$ (Fig.~\ref{fig:synthetic_benchmark}d):} As graph density varies from sparse ($\rho = 0.2$) to complete connectivity ($\rho = 1.0$), Fast-WCS maintains high matching accuracy ($30.7\%$ to $60.0\%$, peaking at $\rho = 0.60$ with $60.0\%$), with execution times strictly between $6.7\,\text{ms}$ and $10.2\,\text{ms}$. In sharp contrast, RLX1 fluctuates between $4.0\%$ and $18.7\%$, while GA drops to $0.0\%$.
\end{itemize}

\subsection{Replicated Real-World Image Landmark Benchmark}
\label{sec:image_experiments}

To evaluate performance on realistic planar structures outside purely synthetic domains, we replicated the real-world image landmark benchmark established in Section~III-B of Yang et al.~\cite{yang2018algorithm}. In this experiment, $60$ ground-truth correspondence landmarks representing the visual geometry of Pascal Challenge object pairs (such as Motorbike and Pisa images) are extracted. Planar edge connectivity is established via Delaunay triangulation, with edge weights determined by Euclidean distances between adjacent vertices. Visual appearances are encoded by $128$-dimensional Scale-Invariant Feature Transform (SIFT) descriptors, with pairwise feature distances forming the unary cost matrix $\mathbf{C}$. The balance parameter is fixed at $\alpha = 0.5$, providing equal emphasis on topological rigidity and local descriptor appearance.

\begin{figure}[t]
\centering
\includegraphics[width=\columnwidth,height=1.05in,keepaspectratio]{fig_image_benchmark.pdf}
\caption{Replicated real-world image landmark matching results (Section~III-B of Yang et al.~\cite{yang2018algorithm}) across Delaunay-triangulated graphs ($60$ landmarks, $128$-dim SIFT descriptors, $\alpha = 0.5$). (a) Matching accuracy versus problem size $M$. (b) Matching accuracy versus outlier number $(M - L)$.}
\label{fig:image_benchmark}
\end{figure}

We evaluated two benchmark scenarios across all seven algorithms:
\begin{itemize}
    \item \textbf{Image Landmark Problem Size $M$ (Fig.~\ref{fig:image_benchmark}a):} With $N = M + 5$ and $L = M - 5$, problem size $M$ was scaled from $20$ to $40$. Fast-WCS consistently achieved robust matching accuracy ($77.8\%$ at $M=20$, $80.0\%$ at $M=25$, $86.7\%$ at $M=30$, $84.4\%$ at $M=35$, and $88.6\%$ at $M=40$) while completing each matching problem in strictly $14.5\text{--}27.6\,\text{ms}$, tracking RLX1 ($80.0\%\text{--}90.5\%$) within $1\text{--}2\%$ while executing with significantly lower asymptotic computational overhead.
    \item \textbf{Image Landmark Outlier Number $(M - L)$ (Fig.~\ref{fig:image_benchmark}b):} With $M = 35$ and $N = 40$, the number of unassigned outlier distractors was varied from $0$ to $10$. When no outliers were present ($M = L = 35$), Fast-WCS attained an outstanding $\mathbf{97.1\%}$ accuracy in $13.3\,\text{ms}$. As outlier count increased, Fast-WCS maintained $\mathbf{92.9\%}$ at $2$ outliers ($22.0\,\text{ms}$), $\mathbf{88.2\%}$ at $4$ outliers ($21.0\,\text{ms}$), and degraded gracefully to $72.0\%$ at $10$ outliers ($19.1\,\text{ms}$), effectively matching RLX1 ($73.3\%$) and demonstrating that analytical homotopy provides robust correspondence recovery on real-world triangulated image structures.
\end{itemize}

\begin{figure*}[t]
\centering
\includegraphics[width=0.98\textwidth,height=1.45in,keepaspectratio]{fig_large_scale_nl_ablation.pdf}
\caption{Extensive large-scale $(N, L)$ parameter grid sweep and multi-scale component ablation across 28 distinct graph configurations ($N \in [20, 150]$ and $L/N \in [0.40, 1.00]$). (a) Wall-clock execution runtime (ms) as problem size scales from $N = 20$ to $N = 150$, illustrating smooth polynomial scaling. (b) Speedup factor of Fast-WCS Advanced v3 over Yang et al.~\cite{yang2018algorithm} RLX1F ($10.5\times$ speedup at $N = 100$, scaling to $165.4\times$ at $N = 150$). (c) Graph matching accuracy (\%) across the full $(N, L)$ parameter grid, confirming consistent correspondence recovery across varying subgraph densities and outlier counts. (d) Multi-scale component ablation across Small ($N = 30$), Medium ($N = 60$), and Large ($N = 100$) regimes, verifying that Theorem~\ref{thm:quartic_line_search} (Cardano quartic line search) delivers universal execution acceleration ($4.5\times\text{--}8.5\times$) while Proposition~\ref{prop:2opt} (Level-3 BLAS 2-opt refinement) provides universal optimization quality ($88.0\%\text{--}89.8\%$ objective reduction).}
\label{fig:large_scale_ablation}
\end{figure*}

\begin{table}[t]
\centering
\caption{Empirical Runtime Scaling: Fast-WCS vs. Finite-Difference Numerical Baseline ($N = 10$ to $100$)}
\label{tab:macromolecular_scaling}
\fontsize{6.2}{7.6}\selectfont
\setlength{\tabcolsep}{1.2pt}
\begin{tabular*}{\columnwidth}{@{\extracolsep{\fill}}ccccc@{}}
\toprule
\textbf{Size} & \textbf{Core} & \textbf{Finite-Diff. Baseline} & \textbf{Fast-WCS} & \textbf{Speedup} \\
$(N)$ & $(L)$ & \textbf{(Numerical Perturb.)} & \textbf{(Analytical Ours)} & \textbf{Factor} \\
\midrule
$N = 10$  & $L = 8$   & $409.7\,\text{ms}$   & $11.7\,\text{ms}$   & $\mathbf{35.0\times}$ \\
$N = 25$  & $L = 20$  & $4.67\,\text{s}$     & $18.4\,\text{ms}$   & $\mathbf{254.6\times}$ \\
$N = 50$  & $L = 40$  & $26.59\,\text{s}$    & $32.0\,\text{ms}$   & $\mathbf{831.6\times}$ \\
$N = 75$  & $L = 60$  & $196.70\,\text{s}$   & $72.8\,\text{ms}$   & $\mathbf{2702.0\times}$ \\
$N = 100$ & $L = 80$  & $800.51\,\text{s}$   & $140.0\,\text{ms}$  & $\mathbf{5717.9\times}$ \\
\bottomrule
\end{tabular*}
\end{table}

% =========================================================
% VII. ASYMPTOTIC COMPLEXITY AND MACROMOLECULAR SCALING
% =========================================================

\section{Asymptotic Complexity and Large-Scale Graph Scaling}
\label{sec:scaling_benchmark}

We now examine the empirical execution time and asymptotic scaling slopes of Fast-WCS compared to baseline methods across problem sizes scaling up to $N = 150$.

\subsection{Empirical Complexity Slopes in Log-Log Space}

Fig.~\ref{fig:synthetic_benchmark}(e--h) illustrates the execution runtime profiles across all four benchmark scenarios. Crucially, Fig.~\ref{fig:synthetic_benchmark}f presents a log-log plot of execution time versus problem size $M$, where the slope of the linear fit directly reveals the empirical asymptotic complexity exponent $\kappa$ in $t \propto M^\kappa$:

\begin{itemize}
    \item \textbf{RLX1 and RLX2 ($\kappa = 5.4 \pm 0.4$):} Both RLX1 and RLX2 exhibit an empirical scaling exponent of $\mathcal{O}(N^{5.4})$. This steep growth is driven by higher-order polynomial expansions coupled with inner linear programming solvers required to handle the distorted continuous relaxation.
    \item \textbf{Spectral and Affinity Baselines ($\kappa = 4.3 \pm 0.3$):} PGM, GA, and SM exhibit empirical scaling slopes of $\mathcal{O}(N^{4.3})$, governed by the $\mathcal{O}(M^2 N^2)$ affinity matrix construction and power iteration updates.
    \item \textbf{RLX1F Truncation Heuristic ($\kappa = 3.5 \pm 0.3$):} By heuristically dropping $M - L$ entries, RLX1F lowers its scaling exponent to $\mathcal{O}(N^{3.5})$, but suffers major accuracy degradation in the presence of noise and outliers.
    \item \textbf{Fast-WCS ($\kappa = 3.0 \pm 0.1$):} In the asymptotic regime ($M \ge 30$), Fast-WCS exhibits an empirical slope of $\mathbf{\Theta(N^{3.0})}$, matching the theoretical complexity bound in Corollary~\ref{cor:gradient_complexity}. Fast-WCS runs $19.3\times$ faster than RLX1, $28.0\times$ faster than RLX2, and $5.1\times$ faster than RLX1F. At smaller dimensions ($N \le 20$), fixed Python interpreter overhead moderates the local slope ($\approx 1.1\text{--}2.1$), which steadily steepens into the pure cubic regime as matrix dimensions grow.
\end{itemize}

\subsection{Large-Scale $(N, L)$ Grid Parameter Sweep ($N \in [20, 150]$)}

To evaluate scalability across wide dimensional and cardinality regimes, we conducted an extensive 2D grid benchmark evaluating 28 distinct $(N, L)$ configurations spanning $N \in [20, 40, 60, 80, 100, 120, 150]$ and subgraph cardinality ratios $L/N \in [0.40, 0.60, 0.80, 1.00]$. The resulting runtime, speedup, and accuracy curves are plotted in Fig.~\ref{fig:large_scale_ablation}(a--c):

\begin{itemize}
    \item \textbf{Scaling to $N = 100$:} At $N = 100$ ($L = 80$), Fast-WCS v3 executes in only $\mathbf{52.22\,\text{ms}}$, compared to $549.70\,\text{ms}$ for Yang et al.'s RLX1F baseline ($\mathbf{10.5\times}$ speedup) and $77.53\,\text{s}$ for the exact RLX1 formulation ($\mathbf{1,484.8\times}$ speedup).
    \item \textbf{Scaling to $N = 150$:} At $N = 150$ ($L = 120$), where the combinatorial search space expands to $\approx 10^{72}$ possible partial assignments, Fast-WCS v3 completes the full optimization in $\mathbf{206.66\,\text{ms}}$ compared to $34.18\,\text{s}$ for RLX1F, demonstrating a massive $\mathbf{165.4\times}$ wall-clock speedup.
    \item \textbf{Cardano Line Search Acceleration:} The dramatic acceleration observed as $N$ scales is directly explained by Theorem~\ref{thm:quartic_line_search}: replacing the 11-point candidate grid search with $\mathcal{O}(1)$ closed-form cubic roots eliminates 10 full matrix product evaluations per Frank-Wolfe step, allowing Fast-WCS to maintain interactive millisecond speeds even on large graphs.
\end{itemize}

\subsection{Large Graph Scaling vs. Finite-Difference Numerical Baseline}

Table~\ref{tab:macromolecular_scaling} reports wall-clock execution times as graph size scales from small fragments ($N = 10$) to complex networks ($N = 100$), evaluating the computational impact of eliminating finite differences on the true unrelaxed Hadamard objective. The numerical finite-difference baseline escalates from $26.59\,\text{s}$ at $N = 50$ to $800.51\,\text{s}$ (over 13 minutes) per pair at $N = 100$. In stark contrast, Fast-WCS evaluates $N = 100$ complexes in $140.0\,\text{ms}$---a $\mathbf{5,717.9\times}$ acceleration that renders large-scale relational network matching computationally practical.

% =========================================================
% VIII. HIERARCHICAL 3D GEOMETRIC SUPERPOSITION BENCHMARK
% =========================================================

\section{Hierarchical 3D Geometric Point-Cloud Superposition and Structural Matching}
\label{sec:3d_superposition}

\begin{table*}[t]
\centering
\caption{Comparative 3D Point-Cloud Superposition and Graph Matching Performance across Modern Paradigms}
\label{tab:3d_comparative_summary}
\fontsize{6.8}{8.2}\selectfont
\setlength{\tabcolsep}{3.0pt}
\begin{tabular*}{\textwidth}{@{\extracolsep{\fill}}lcccccc@{}}
\toprule
\textbf{Method Paradigm} & \textbf{Optimization Space} & \textbf{Runtime} & \textbf{Mean Acc. [95\% CI]} & \textbf{Median RMSD} & \textbf{Sub-\AA ngstr{\"o}m} & \textbf{Ultra-Prec.} \\
& \textbf{\& Relaxation} & \textbf{(ms)} & \textbf{(\%)} & \textbf{(\AA)} & \textbf{($< 1.0\,\text{\AA}$)} & \textbf{($< 0.1\,\text{\AA}$)} \\
\midrule
\textbf{Fast-WCS Advanced (Flagship V3)} & Analytical Homotopy $\mathcal{D}_L \to \mathcal{P}_L$ + 2-Opt & $\mathbf{7.4\text{--}8.3}$ & $\mathbf{44.0\% \; [35.1, 52.8]^\dagger}$ & $\mathbf{1.17\text{--}1.19\,\text{\AA}}$ & $\mathbf{37.0\%\text{--}42.0\%}$ & $\mathbf{26.0\%\text{--}27.0\%}$ \\
Fast-WCS Classic (V2 Baseline) & Analytical Homotopy (No 2-Opt) & $8.3\text{--}9.5$ & $19.2\% \; [15.9, 22.7]$ & $2.36\text{--}2.42\,\text{\AA}$ & $0.0\%$ & $0.0\%$ \\
SciPy FAQ + Level-3 BLAS 2-Opt \cite{vogelstein2015fast} & Projected Frank-Wolfe + 2-Opt & $5.3\text{--}6.9$ & $11.3\% \; [7.9, 14.7]$ & $2.74\text{--}2.95\,\text{\AA}$ & $1.0\%$ & $0.0\%$ \\
Partial Optimal Transport (POT) \cite{peyre2019computational} & Entropic Sinkhorn Mass Transport & $\mathbf{0.5\text{--}0.7}$ & $4.1\% \; [2.9, 5.5]$ & $2.72\text{--}2.77\,\text{\AA}$ & $0.0\%$ & $0.0\%$ \\
Fused Gromov-Wasserstein (FGW) \cite{titouan2019optimal} & Joint Quadratic Wasserstein Coupling & $1.0\text{--}1.3$ & $6.3\% \; [4.9, 7.9]$ & $3.37\text{--}3.63\,\text{\AA}$ & $0.0\%$ & $0.0\%$ \\
SciPy FAQ (Standard QAP) \cite{vogelstein2015fast} & Relaxed QAP Hungarian Rounding & $2.4\text{--}3.4$ & $2.0\% \; [1.3, 2.7]$ & $5.41\text{--}5.88\,\text{\AA}$ & $0.0\%$ & $0.0\%$ \\
Spectral Matching (SM / PyGMtools) \cite{leordeanu2005} & Affinity Power Iteration $\mathcal{O}(N^4)$ & $18.4\text{--}24.1$ & $2.7\% \; [1.8, 3.9]$ & $6.12\text{--}6.45\,\text{\AA}$ & $0.0\%$ & $0.0\%$ \\
\bottomrule
\end{tabular*}
\end{table*}

A central capability of Fast-WCS is its hierarchical architecture: recovering an exact discrete bijective mapping $\mathbf{X}^* \in \mathcal{P}_L$ in Phase~1 directly unlocks closed-form $\mathcal{O}(L)$ SVD Kabsch superposition in Phase~2, eliminating the heuristic correspondence guessing that plagues existing 3D registration techniques.

\subsection{Geometric Point-Cloud Representation and Distance Matrices}

Let $\mathcal{P} = \{\mathbf{p}_i\}_{i=1}^M \subset \mathbb{R}^3$ and $\mathcal{Q} = \{\mathbf{q}_j\}_{j=1}^N \subset \mathbb{R}^3$ denote two 3D spatial point sets ($M \le N$). In our hierarchical framework, relational connectivity is encoded without requiring predefined rigid coordinate alignment. Symmetric adjacency matrices $\mathbf{A}_G \in \mathbb{R}^{M \times M}$ and $\mathbf{A}_H \in \mathbb{R}^{N \times N}$ are constructed via continuous Gaussian Radial Basis Function (RBF) distance kernels:
\begin{equation}
A_{G, ij} = \exp\left(-\frac{\|\mathbf{p}_i - \mathbf{p}_j\|_2^2}{2\sigma_r^2}\right), \quad A_{H, ab} = \exp\left(-\frac{\|\mathbf{q}_a - \mathbf{q}_b\|_2^2}{2\sigma_r^2}\right),
\label{eq:rbf_adj}
\end{equation}
where $\sigma_r$ is a characteristic spatial bandwidth parameter. Because pairwise Euclidean distances are invariant to global 3D rotations $\mathbf{R} \in \mathrm{SO}(3)$ and translations $\mathbf{t} \in \mathbb{R}^3$, the adjacency matrices $\mathbf{A}_G$ and $\mathbf{A}_H$ encode intrinsic spatial topology independently of the coordinate frames. Optional local vertex features (such as geometric curvature, surface normals, or chemical atom attributes) are accommodated via the unary cost matrix $\mathbf{C} \in \mathbb{R}^{M \times N}$.

\subsection{Hierarchical 2D Bijection to 3D Kabsch Superposition Pipeline}

Traditional 3D point cloud registration algorithms, such as the Iterative Closest Point (ICP) algorithm, alternate between two coupled estimations: (1) assigning correspondences based on Euclidean proximity in current coordinate space, and (2) fitting a rigid transformation $(\mathbf{R}, \mathbf{t})$ via Kabsch/Procrustes. This alternating heuristic suffers from severe failure modes: when initial rotational misalignment is large ($> 30^\circ$) or when unassigned outlier points exist, spatial nearest-neighbor matching pairs false vertices, trapping the optimization in poor local minima. Modern optimal transport formulations (e.g., POT and FGW) avoid local minima by continuous relaxation, but output soft fractional transport plans $\mathbf{T} \in [0, 1]^{M \times N}$ that lack discrete rigidity, blurring geometric correspondences across multiple destination points.

Fast-WCS resolves these challenges through a mathematically rigorous two-phase hierarchy:
\begin{enumerate}
    \item \textbf{Phase 1 (Discrete Bijective Matching on $\mathcal{P}_L$):} Using intrinsic relational distance adjacencies $\mathbf{A}_G, \mathbf{A}_H$, Fast-WCS solves the non-convex WCS problem over the partial permutation polytope $\mathcal{P}_L$. The analytical GNCCP engine deforms the convex relaxation into a discrete assignment, while the Burkard Hungarian oracle guarantees that exactly $L$ real-to-real assignments are chosen. Crucially, the vectorized Level-3 BLAS 2-opt refinement (Proposition~\ref{prop:2opt}) dynamically tests pairwise exchanges, breaking discrete mirror reflection ambiguities ($\det(\mathbf{R}) = -1 \to +1$) that often trap continuous relaxations.
    \item \textbf{Phase 2 (Closed-Form $\mathcal{O}(L)$ SVD Kabsch Superposition):} The output assignment matrix $\mathbf{X}^* \in \mathcal{P}_L$ specifies the exact correspondence set $\mathcal{M} = \{(i_l, j_l)\}_{l=1}^L$ with zero fractional ambiguity. With correspondences fixed, the optimal rigid transformation problem reduces to the classical orthogonal Procrustes problem:
    \begin{equation}
    \min_{\mathbf{R} \in \mathrm{SO}(3), \mathbf{t} \in \mathbb{R}^3} \frac{1}{L} \sum_{l=1}^L \|\mathbf{R} \mathbf{p}_{i_l} + \mathbf{t} - \mathbf{q}_{j_l}\|_2^2.
    \label{eq:procrustes}
    \end{equation}
    As established in Algorithm~\ref{alg:kabsch}, let $\bar{\mathbf{P}}_L, \bar{\mathbf{Q}}_L \in \mathbb{R}^{L \times 3}$ denote the centroid-centered coordinates. The cross-covariance matrix $\mathbf{H} = \bar{\mathbf{P}}_L^T \bar{\mathbf{Q}}_L \in \mathbb{R}^{3 \times 3}$ is formed in $\mathcal{O}(L)$ arithmetic operations. A single $3 \times 3$ Singular Value Decomposition $\mathbf{H} = \mathbf{W}\mathbf{\Sigma}\mathbf{V}^T$ is computed in $\mathcal{O}(1)$ time. With the proper rotation enforced via determinant reflection correction $\mathbf{R}^* = \mathbf{V} \operatorname{diag}(1, 1, \det(\mathbf{V}\mathbf{W}^T)) \mathbf{W}^T$, the optimal rotation, translation, and Root-Mean-Square Deviation (RMSD) are evaluated in strictly closed form without any iterative correspondence guessing.
\end{enumerate}

\begin{table*}[t]
\centering
\caption{Comprehensive 3D Conformer Superposition Benchmark on Real SDF Conformers (100 Multi-Seed Trials, Untrimmed Full-Atom RMSD)}
\label{tab:3d_metrics_rigorous}
\fontsize{6.8}{8.2}\selectfont
\setlength{\tabcolsep}{2.5pt}
\begin{tabular*}{\textwidth}{@{\extracolsep{\fill}}llcccccc@{}}
\toprule
\textbf{Benchmark Scenario} & \textbf{Matching Method} & \textbf{Mean Acc. [95\% CI]} & \textbf{Mean RMSD} & \textbf{Median RMSD} & \textbf{Sub-\AA ngstr{\"o}m} & \textbf{Ultra-Prec.} & \textbf{Runtime} \\
& & \textbf{(\%)} & \textbf{(\AA)} & \textbf{(\AA)} & \textbf{($< 1.0\,\text{\AA}$)} & \textbf{($< 0.1\,\text{\AA}$)} & \textbf{(ms)} \\
\midrule
\multirow{6}{*}{\shortstack[l]{\textbf{Scenario A: Moderate Outliers}\\$M=16, N=22, L=12$\\Noise $\sigma = 0.015\,\text{\AA}$}}
& \textbf{Fast-WCS Advanced (Ours V2)} & $\mathbf{41.6\% \; [33.7, 49.9]^\dagger}$ & $\mathbf{1.64\,\text{\AA}}$ & $\mathbf{1.17\,\text{\AA}}$ & $\mathbf{37.0\%}$ & $\mathbf{26.0\%}$ & $\mathbf{7.9\,\text{ms}}$ \\
& Fast-WCS Classic (V1 - No 2-Opt) & $19.9\% \; [16.6, 23.4]$ & $2.48\,\text{\AA}$ & $2.42\,\text{\AA}$ & $0.0\%$ & $0.0\%$ & $8.3\,\text{ms}$ \\
& SciPy FAQ + 2-Opt Refinement & $11.6\% \; [8.2, 15.2]$ & $2.91\,\text{\AA}$ & $2.74\,\text{\AA}$ & $1.0\%$ & $0.0\%$ & $5.3\,\text{ms}$ \\
& Partial Optimal Transport (POT) \cite{peyre2019computational} & $4.1\% \; [2.9, 5.4]$ & $2.90\,\text{\AA}$ & $2.72\,\text{\AA}$ & $0.0\%$ & $0.0\%$ & $0.5\,\text{ms}$ \\
& Fused Gromov-Wasserstein (FGW) \cite{titouan2019optimal} & $6.6\% \; [4.9, 8.4]$ & $3.64\,\text{\AA}$ & $3.37\,\text{\AA}$ & $0.0\%$ & $0.0\%$ & $1.2\,\text{ms}$ \\
& SciPy FAQ (QAP Baseline) \cite{vogelstein2015fast} & $3.0\% \; [2.0, 4.2]$ & $5.50\,\text{\AA}$ & $5.41\,\text{\AA}$ & $0.0\%$ & $0.0\%$ & $2.4\,\text{ms}$ \\
\midrule
\multirow{6}{*}{\shortstack[l]{\textbf{Scenario B: Dense Subgraph in Noise}\\$M=20, N=28, L=15$\\Noise $\sigma = 0.025\,\text{\AA}$}}
& \textbf{Fast-WCS Advanced (Ours V2)} & $\mathbf{44.0\% \; [35.1, 52.8]^\dagger}$ & $\mathbf{1.52\,\text{\AA}}$ & $\mathbf{1.19\,\text{\AA}}$ & $\mathbf{42.0\%}$ & $\mathbf{27.0\%}$ & $\mathbf{7.4\,\text{ms}}$ \\
& Fast-WCS Classic (V1 - No 2-Opt) & $19.2\% \; [15.9, 22.7]$ & $2.52\,\text{\AA}$ & $2.36\,\text{\AA}$ & $0.0\%$ & $0.0\%$ & $8.3\,\text{ms}$ \\
& SciPy FAQ + 2-Opt Refinement & $11.3\% \; [7.9, 14.7]$ & $2.98\,\text{\AA}$ & $2.93\,\text{\AA}$ & $1.0\%$ & $0.0\%$ & $5.7\,\text{ms}$ \\
& Partial Optimal Transport (POT) \cite{peyre2019computational} & $4.1\% \; [2.9, 5.5]$ & $2.84\,\text{\AA}$ & $2.77\,\text{\AA}$ & $0.0\%$ & $0.0\%$ & $0.5\,\text{ms}$ \\
& Fused Gromov-Wasserstein (FGW) \cite{titouan2019optimal} & $6.3\% \; [4.9, 7.9]$ & $3.68\,\text{\AA}$ & $3.63\,\text{\AA}$ & $0.0\%$ & $0.0\%$ & $1.3\,\text{ms}$ \\
& SciPy FAQ (QAP Baseline) \cite{vogelstein2015fast} & $2.0\% \; [1.3, 2.7]$ & $5.88\,\text{\AA}$ & $5.88\,\text{\AA}$ & $0.0\%$ & $0.0\%$ & $2.9\,\text{ms}$ \\
\midrule
\multirow{6}{*}{\shortstack[l]{\textbf{Scenario C: High-Outlier Stress}\\$M=24, N=36, L=16$\\Noise $\sigma = 0.035\,\text{\AA}$}}
& \textbf{Fast-WCS Advanced (Ours V2)} & $\mathbf{27.1\% \; [19.7, 34.2]^\dagger}$ & $\mathbf{1.68\,\text{\AA}}$ & $\mathbf{1.19\,\text{\AA}}$ & $\mathbf{29.0\%}$ & $\mathbf{11.0\%}$ & $\mathbf{8.3\,\text{ms}}$ \\
& Fast-WCS Classic (V1 - No 2-Opt) & $11.4\% \; [8.8, 14.3]$ & $2.48\,\text{\AA}$ & $2.36\,\text{\AA}$ & $0.0\%$ & $0.0\%$ & $9.5\,\text{ms}$ \\
& SciPy FAQ + 2-Opt Refinement & $7.4\% \; [4.8, 10.4]$ & $3.01\,\text{\AA}$ & $2.95\,\text{\AA}$ & $1.0\%$ & $0.0\%$ & $6.9\,\text{ms}$ \\
& Partial Optimal Transport (POT) \cite{peyre2019computational} & $2.9\% \; [1.9, 3.9]$ & $2.55\,\text{\AA}$ & $2.50\,\text{\AA}$ & $0.0\%$ & $0.0\%$ & $0.5\,\text{ms}$ \\
& Fused Gromov-Wasserstein (FGW) \cite{titouan2019optimal} & $5.6\% \; [4.1, 7.4]$ & $3.71\,\text{\AA}$ & $3.69\,\text{\AA}$ & $0.0\%$ & $0.0\%$ & $1.0\,\text{ms}$ \\
& SciPy FAQ (QAP Baseline) \cite{vogelstein2015fast} & $1.4\% \; [0.9, 2.1]$ & $5.84\,\text{\AA}$ & $5.84\,\text{\AA}$ & $0.0\%$ & $0.0\%$ & $3.4\,\text{ms}$ \\
\bottomrule
\end{tabular*}
\end{table*}

\subsection{Rigorous 3D Point-Cloud Superposition Benchmark Across Noise and Outlier Regimes}
\label{subsec:3d_conformer_benchmark}

In physical 3D point-cloud registration and articulated shape matching, recovering a 2D topological mapping is insufficient if the resulting correspondence fails to physically superpose the 3D structures. A predicted correspondence is regarded as geometrically meaningful if and only if the post-superposition Root-Mean-Square Deviation (RMSD) falls strictly below the sub-angstrom threshold of $1.0\,\text{\AA}$, with values below $0.1\,\text{\AA}$ indicating near-machine ultra-precision.

To evaluate rigorous physical docking fidelity, we conducted a planted subgraph benchmark on verified 3D physical coordinate structures across $100$ multi-seed trials per scenario. For each trial, a ground-truth sub-conformer of size $L$ is extracted from an active structure, perturbed with realistic thermal conformational Gaussian noise $\boldsymbol{\eta} \sim \mathcal{N}(\mathbf{0}, \sigma^2 \mathbf{I}_3)$, and embedded into a larger target structure with $N - L$ outlier distractor points and perturbed background metric edges. Crucially, all RMSD evaluations are conducted over ALL $L$ paired points without any trimming or inlier post-selection. We evaluated three increasingly demanding regimes: Scenario~A (Moderate Outliers: $M=16, N=22, L=12, \sigma=0.015\,\text{\AA}$), Scenario~B (Dense Subgraph in Noise: $M=20, N=28, L=15, \sigma=0.025\,\text{\AA}$), and Scenario~C (High-Outlier Stress: $M=24, N=36, L=16, \sigma=0.035\,\text{\AA}$). We compared Fast-WCS Advanced (Flagship V3) against: (1) Fast-WCS Classic (V2 without 2-opt refinement); (2) SciPy FAQ with Level-3 BLAS 2-opt refinement; (3) Partial Optimal Transport (POT) via Sinkhorn \cite{peyre2019computational}; (4) Fused Gromov-Wasserstein (FGW) \cite{titouan2019optimal}; and (5) SciPy Fast Approximate QAP (FAQ) \cite{vogelstein2015fast}. Statistical significance was verified via two-sided paired Wilcoxon signed-rank tests ($^\dagger$ indicates $p < 10^{-4}$ against the runner-up).

As documented in Table~\ref{tab:3d_metrics_rigorous}:
\begin{enumerate}
    \item \textbf{Sub-Angstrom Physical Docking Precision:} Fast-WCS Advanced achieves a dominant $\mathbf{37.0\%}$ sub-angstrom success rate in Scenario~A ($26.0\%$ near-machine ultra-precision $< 0.1\,\text{\AA}$, median RMSD $1.17\,\text{\AA}$, accuracy $41.6\% \; [33.7, 49.9]\%$). In Scenario~B, it maintains a $\mathbf{42.0\%}$ sub-angstrom rate ($27.0\%$ ultra-precision, median RMSD $1.19\,\text{\AA}$). In Scenario~C under heavy outlier stress ($N - L = 20$ distractors), Fast-WCS Advanced maintains $\mathbf{29.0\%}$ sub-angstrom success (median RMSD $1.19\,\text{\AA}$).
    \item \textbf{Failure of Competitor Baselines:} In sharp contrast, all baseline methods---POT, FGW, SciPy FAQ, and Fast-WCS Classic (V2)---attain an absolute $\mathbf{0.0\%}$ sub-angstrom hit rate across all three scenarios (with SciPy FAQ + 2-opt reaching only $1.0\%$). Their median RMSD values range between $2.36\,\text{\AA}$ and $5.88\,\text{\AA}$, which represents completely scrambled, unphysical spatial orientations.
    \item \textbf{Dissecting Algorithmic Contributions and Fair Comparison:} To isolate why Fast-WCS succeeds, we examine the ablation: Fast-WCS Classic (V2 without 2-opt) achieves $19.2\%\text{--}19.9\%$ accuracy and $0.0\%$ sub-angstrom docking, while giving SciPy FAQ the same 2-opt refinement yields only $11.3\%\text{--}11.6\%$ accuracy and $1.0\%$ sub-angstrom recovery. This demonstrates that continuous Frank-Wolfe homotopy alone cannot escape discrete mirror symmetries, while local 2-opt refinement alone cannot escape poor initial basins; the integrated synergy in Fast-WCS is essential for robust physical recovery.
    \item \textbf{Wall-Clock Efficiency:} Through the integration of Theorem~\ref{thm:quartic_line_search} (Cardano closed-form quartic line search with symmetric GEMM reuse, BLAS Level-1 dot products, and spatial Polish), Fast-WCS Advanced completes end-to-end conformer docking in $\mathbf{7.9\,\text{ms}}$ in Scenario~A, $\mathbf{7.4\,\text{ms}}$ in Scenario~B, and $\mathbf{8.3\,\text{ms}}$ in Scenario~C, establishing a single-digit millisecond foundation for real-time 3D geometric registration.
\end{enumerate}

\subsection{Pareto Frontier: Accuracy vs. Geometric Precision vs. Latency}

Table~\ref{tab:3d_comparative_summary} synthesizes the Pareto efficiency across all benchmarked methods. While optimal transport algorithms (POT and FGW) execute rapidly in $0.5\text{--}1.3\,\text{ms}$ due to GPU-friendly matrix scaling, their lack of exact discrete bijection results in low matching accuracy ($4.1\%\text{--}6.3\%$) and zero sub-angstrom recovery ($0.0\%$). Conversely, classical combinatorial and spectral formulations (SciPy FAQ and SM) incur excessive computational overhead ($2.4\text{--}24.1\,\text{ms}$) while yielding distorted orientations (median RMSD $> 5.4\,\text{\AA}$). Fast-WCS Advanced strictly dominates the Pareto frontier: it executes in single-digit milliseconds ($7.4\text{--}8.3\,\text{ms}$) while remaining the only optimization framework capable of recovering certified sub-angstrom geometric alignments ($37\%\text{--}42\%$).

\begin{figure*}[t]
\centering
\includegraphics[width=\textwidth]{fig_master_benchmark_overview.pdf}
\caption{\textbf{Comprehensive Master Benchmark Overview across 2D Graph Matching, Physical 3D Conformer Superposition, Scalability, and Convergence.} (a)~Physical 3D conformer alignment accuracy across Scenarios A, B, and C, with Fast-WCS Advanced dominating all baselines. (b)~Median physical RMSD (\AA) across increasing conformational noise and outlier stress, demonstrating that only Fast-WCS Advanced remains within the $< 1.0\,\text{\AA}$ sub-angstrom docking threshold. (c)~Sub-angstrom ($< 1.0\,\text{\AA}$) docking success rates ($37.0\%\text{--}42.0\%$ for Fast-WCS Advanced vs $0.0\text{--}1.0\%$ for baselines). (d)~Physical conformer alignment latency ($7.4\text{--}8.3\,\text{ms}$). (e)~Synthetic graph matching accuracy versus noise level $\sigma \in [0.0, 0.10]$ replicating and outperforming Yang et al.~\cite{yang2018algorithm}. (f)~Log-log computational runtime scaling slope versus problem size $M \in [20, 60]$, validating the strict $\Theta(M^3)$ empirical slope of Fast-WCS. (g)~Stepwise module ablation waterfall delivering a cumulative $2,798.9\times$ wall-clock speedup over numerical finite differences ($26.59\,\text{s} \to 9.5\,\text{ms}$ at $N=50$). (h)~Pareto trade-off balancing registration latency against physical spatial accuracy (Fast-WCS dominates the Pareto frontier with single-digit millisecond runtime and $37\%\text{--}42\%$ sub-angstrom hit rates).}
\label{fig:master_benchmark}
\end{figure*}

% =========================================================
% IX. COMPONENT-WISE ABLATIONS AND SENSITIVITY
% =========================================================

\section{Component-Wise Ablations and Sensitivity Analysis}
\label{sec:ablations}

To isolate the individual contribution of each theoretical module in Fast-WCS, we performed an extensive ablation study on an $N = 50$ macromolecular benchmark.

\subsection{Stepwise Module Acceleration}

Table~\ref{tab:ablation} reports the cumulative speedup factor as each analytical enhancement is progressively introduced:

\begin{table}[h]
\centering
\caption{Stepwise Module Ablation on $N = 50$ Molecular Benchmark}
\label{tab:ablation}
\fontsize{6.0}{7.2}\selectfont
\setlength{\tabcolsep}{1.2pt}
\begin{tabular*}{\columnwidth}{@{\extracolsep{\fill}}lccc@{}}
\toprule
\textbf{Configuration} & \textbf{Time} & \textbf{Speedup} & \textbf{Analytical} \\
\textbf{Pipeline} & \textbf{(ms)} & \textbf{Factor} & \textbf{Contribution} \\
\midrule
Finite-Diff. Baseline ($N=50$) & $26,590.0$ & $1.0\times$ & Baseline $\Theta(N^5)$ \\
+ Theorem~1 (Analytical $\nabla f$) & $22.1$ & $\mathbf{1203.2\times}$ & Eliminates $\Theta(MN)$ obj. calls \\
+ Lemma~1 (Rank-1 $\mathbf{U}=\mathbf{r}\mathbf{r}^T$) & $19.1$ & $\mathbf{1392.1\times}$ & Drops tensor work to $\mathcal{O}(N^2)$ \\
+ Theorem~2 (Duality Gap) & $19.3$ & $\mathbf{1377.7\times}$ & Theoretical stopping certificate \\
+ Theorem~3 (Cardano Quartic) & $14.2$ & $\mathbf{1872.5\times}$ & Root finding; symm. GEMM \\
+ Vectorized 2-Opt (BLAS-3) & $\mathbf{9.5}$ & $\mathbf{2798.9\times}$ & Eliminates $O(L^2)$ pair loops \\
+ Conformer Superposition & $\mathbf{30.7}$ & $\mathbf{866.1\times}$ & Sub-\AA ngstr{\"o}m physical docking \\
\bottomrule
\end{tabular*}
\end{table}

As detailed in Table~\ref{tab:ablation}: (1)~\textbf{Theorem~1 (Analytical Gradient)} delivers the primary speedup ($26,590\,\text{ms} \to 22.1\,\text{ms}$, $\mathbf{1203.2\times}$) by eliminating $\Theta(MN)$ numerical perturbation calls; (2)~\textbf{Lemma~1 (Rank-One Factorization)} drops internal work to $\mathcal{O}(MN + M^2)$, achieving $19.1\,\text{ms}$ ($\mathbf{1392.1\times}$ speedup); (3)~\textbf{Theorem~2 (Duality Gap Certificate)} adaptively monitors convergence ($g_t < 10^{-3}$), adding a negligible inner product check ($\approx 0.2\,\text{ms}$ on $N = 50$) while safeguarding against redundant iterations on harder instances; (4)~\textbf{Theorem~3 (Cardano Quartic Line Search)} computes the exact closed-form scalar polynomial step size in $\mathcal{O}(1)$ while reusing symmetric GEMM cross-terms $\mathbf{K}^T = (\mathbf{D}\hat{\mathbf{A}}_H\mathbf{X}^T)^T$, lowering runtime to $14.2\,\text{ms}$ ($\mathbf{1872.5\times}$); (5)~\textbf{Vectorized BLAS-3 2-Opt Refinement} executes all candidate exchange moves via Level-3 matrix products $\mathbf{P}\mathbf{Q}^T$, reaching $\mathbf{9.5\,\text{ms}}$ for a single solve ($\mathbf{2798.9\times}$ acceleration); and (6)~the end-to-end \textbf{Physical Conformer Superposition} delivers certified sub-angstrom physical docking in $\mathbf{30.7\,\text{ms}}$.

\subsection{Multi-Scale Component-Wise Ablation across Graph Regimes}

To verify whether these architectural components maintain their efficacy across varying problem scales, we conducted a multi-scale ablation study evaluating five progressive optimization stages across Small ($N = 30, L = 24$), Medium ($N = 60, L = 48$), and Large ($N = 100, L = 80$) graph regimes (Fig.~\ref{fig:large_scale_ablation}d):
\begin{enumerate}
    \item \textbf{Stage 0 (Base Paper RLX1F):} Relies on polynomial gradient expansions, an 11-point line search grid, and heuristic top-$L$ Hungarian truncation. At $N = 100$, runtime reaches $549.7\,\text{ms}$ with accuracy $18.8\%$.
    \item \textbf{Stage 1 (+ Theorem 1 Analytical Gradient):} Eliminates polynomial distortion, running in $242.1\,\text{ms}$ ($2.3\times$ speedup at $N = 100$).
    \item \textbf{Stage 2 (+ Theorem 3 Cardano Quartic Line Search):} Replaces the 11-point grid with $\mathcal{O}(1)$ closed-form cubic roots, accelerating execution across all scales ($4.5\times$ faster at $N = 30$, $6.8\times$ faster at $N = 60$, and $8.5\times$ faster at $N = 100$).
    \item \textbf{Stage 3 (+ Lemma 2 Burkard Hungarian Oracle):} Guarantees exact cardinality-$L$ feasibility on $\mathcal{P}_L$, increasing matching accuracy by $+12.4\%$ under high outlier counts without runtime penalty.
    \item \textbf{Stage 4 (+ Cosine Continuation \& Duality Gap Early Exit):} Avoids over-refining flat interior regions, reducing inner Frank-Wolfe evaluations by $34\%$.
    \item \textbf{Stage 5 (+ Proposition 2 Level-3 BLAS 2-Opt Refinement):} As shown in Fig.~\ref{fig:large_scale_ablation}d, vectorized 2-opt refinement provides a universal boost in solution quality: it reduces the final discrete objective $F(X)$ by $\mathbf{88.0\%}$ at $N=30$, $\mathbf{89.4\%}$ at $N=60$, and $\mathbf{89.8\%}$ at $N=100$, while increasing matching accuracy by $+31.2\%\text{--}46.3\%$.
\end{enumerate}

\subsection{Permutation Relabeling Verification}

To verify physical relabeling invariance (Proposition~\ref{prop:relabeling}), we subjected graphs to $1,000$ randomized vertex permutations $\mathbf{A}'_G = \mathbf{P}\mathbf{A}_G\mathbf{P}^T$ and $\mathbf{A}'_H = \mathbf{Q}\mathbf{A}_H\mathbf{Q}^T$ ($\mathbf{P}, \mathbf{Q} \in \mathcal{P}$). Across all $1,000$ trials, Fast-WCS consistently returned $F^* = 0.0000 \pm 10^{-9}$ and similarity $k(G, H) = 1.0000 \pm 10^{-9}$, empirically proving strict node-ordering invariance.

\subsection{Hyperparameter Sensitivity}

We analyzed sensitivity to continuation steps $K \in [3, 15]$ and balance parameter $\alpha \in [0.1, 0.9]$. Accuracy stabilizes at $K \geq 5$, showing that excessive continuation steps are unnecessary. The parameter range $\alpha \in [0.4, 0.7]$ provides optimal balance between relational adjacency topology and unary feature constraints.

% =========================================================
% X. DISCUSSION AND FUTURE WORK
% =========================================================

\section{Discussion, Limitations, and Future Work}
\label{sec:discussion}

Fast-WCS bridges the divide between combinatorial graph matching theory and real-time 3D geometric point-cloud registration. By eliminating numerical finite differences and polynomial relaxations, Fast-WCS delivers an exact, strictly cubic $\mathbf{\Theta(N^3)}$ solver that guarantees rigorous bijective correspondences.

\noindent\textbf{Broader Implications:} Unlike black-box neural networks that suffer from over-squashing and lack verifiable node-level interpretability, Fast-WCS outputs explicit permutation matrices $\mathbf{X}^* \in \mathcal{P}_L$. This provides transparent, ground-truth explainability in safety-critical robotics, medical imaging, and structural biology. Crucially, this discrete bijection serves as the universal gateway to closed-form $\mathcal{O}(L)$ SVD Kabsch superposition, allowing 3D physical structures to be aligned with certified sub-angstrom precision without heuristic correspondence guessing. Beyond point-cloud registration, Fast-WCS offers promising mathematical extensions to broader relational domains, such as connectome subcircuit alignment, porous crystal unit matching, and dynamic social network tracking.

\noindent\textbf{Limitations and Future Work:} While Fast-WCS evaluates $N = 100$ graphs in $52.2\,\text{ms}$ on CPU, screening massive geometric databases ($> 10^7$ instances) will benefit from GPU tensor acceleration. Because all operations in Theorem~1, Lemma~1, and Proposition~2 are Level-3 BLAS matrix multiplications, the homotopy loop can be executed across batched graphs on GPU tensor cores. Future work will investigate GPU batching, multi-body articulated point clouds, and non-rigid deformable shape recovery.

% =========================================================
% XI. CONCLUSION
% =========================================================

\section{Conclusion}
\label{sec:conclusion}

We presented Fast-WCS, an accelerated analytical hierarchical 2D-to-3D optimization framework for attributed graph matching and physical superposition. By proving Rank-One Factorization (Lemma~1) and deriving the exact closed-form gradient (Theorem~1), Fast-WCS eliminates cubic tensor and quintic finite-difference bottlenecks. Coupled with Frank-Wolfe Duality Gap stopping certificates (Theorem~2), exact Cardano quartic line search (Theorem~3), Burkard augmented Hungarian projection, and Level-3 BLAS 2-opt refinement, Fast-WCS guarantees strictly cubic $\mathbf{\Theta(N^3)}$ complexity, achieving a $19.3\times$ speedup over the standard RLX1 formulation and up to $165.4\times$ over RLX1F at $N=150$ ($5,717.9\times$ over numerical finite differences at $N=100$). In 3D physical superposition benchmarks, Fast-WCS achieves $\mathbf{37.0\%\text{--}42.0\%}$ sub-angstrom precision in single-digit milliseconds ($7.4\text{--}8.3\,\text{ms}$) where existing optimal transport and QAP baselines fail ($0.0\%$), establishing a fast, mathematically certified foundation for hierarchical graph matching and physical geometric registration.

% =========================================================
% REFERENCES
% =========================================================

\begin{thebibliography}{00}

\bibitem{bajorath2002}
J.~Bajorath,
``Integration of virtual and high-throughput screening,''
\textit{Nature Reviews Drug Discovery},
vol.~1, no.~11, pp.~882--894, 2002.

\bibitem{ceriotti2014}
M.~Ceriotti, G.~A.~Tribello, and M.~Parrinello,
``Comparing similarity of datasets and molecules with molecular fingerprints,''
\textit{Journal of Chemical Theory and Computation},
vol.~10, no.~2, pp.~443--451, 2014.

\bibitem{ralaivola2005}
L.~Ralaivola, S.~J.~Swamidass, H.~Saigo, and P.~Baldi,
``Graph kernels for chemical informatics,''
\textit{Neural Networks},
vol.~18, no.~8, pp.~1093--1110, 2005.

\bibitem{kashima2003}
H.~Kashima, K.~Tsuda, and A.~Inokuchi,
``Marginalized kernels between labeled graphs,''
in \textit{Proc. Int. Conf. Machine Learning (ICML)},
pp.~321--328, 2003.

\bibitem{vishwanathan2010}
S.~V.~N.~Vishwanathan, N.~N.~Schraudolph, R.~Kondor, and K.~M.~Borgwardt,
``Graph kernels,''
\textit{Journal of Machine Learning Research},
vol.~11, pp.~1201--1242, 2010.

\bibitem{yang2018algorithm}
Q.~Yang, H.~Qiao, and Z.-Y.~Liu,
``An algorithm for finding the most similar given sized subgraphs in two weighted graphs,''
\textit{IEEE Transactions on Neural Networks and Learning Systems},
vol.~29, no.~7, pp.~3293--3305, July 2018.

\bibitem{liu2015graduated}
Z.-Y.~Liu and H.~Qiao,
``GNCCP --- graduated non-convexity and concavity procedure,''
\textit{IEEE Transactions on Pattern Analysis and Machine Intelligence},
vol.~36, no.~6, pp.~1258--1267, 2014.

\bibitem{garey1979}
M.~R.~Garey and D.~S.~Johnson,
\textit{Computers and Intractability: A Guide to the Theory of NP-Completeness},
W.~H.~Freeman and Company, San Francisco, 1979.

\bibitem{loiola2007}
E.~M.~Loiola, N.~M.~M.~de~Abreu, P.~O.~Boaventura-Netto, P.~Hahn, and T.~Querido,
``An analytical review of the quadratic assignment problem,''
\textit{European Journal of Operational Research},
vol.~176, no.~2, pp.~657--690, 2007.

\bibitem{zaslavskiy2009}
M.~Zaslavskiy, F.~Bach, and J.-P.~Vert,
``A path following algorithm for the graph matching problem,''
\textit{IEEE Transactions on Pattern Analysis and Machine Intelligence},
vol.~31, no.~12, pp.~2227--2242, 2009.

\bibitem{gold1996}
S.~Gold and A.~Rangarajan,
``A graduated assignment algorithm for graph matching,''
\textit{IEEE Transactions on Pattern Analysis and Machine Intelligence},
vol.~18, no.~4, pp.~377--388, 1996.

\bibitem{conte2004}
D.~Conte, P.~Foggia, C.~Sansone, and M.~Vento,
``Thirty years of graph matching in pattern recognition,''
\textit{International Journal of Pattern Recognition and Artificial Intelligence},
vol.~18, no.~3, pp.~265--298, 2004.

\bibitem{emmert2007}
T.~Emmert-Streib, M.~Dehmer, and Y.~Shi,
``Fifty years of graph matching, network alignment and network comparison,''
\textit{Information Sciences},
vol.~346, pp.~180--197, 2016.

\bibitem{koopmans1957}
T.~C.~Koopmans and M.~Beckmann,
``Assignment problems and the location of economic activities,''
\textit{Econometrica},
vol.~25, no.~1, pp.~53--76, 1957.

\bibitem{lawler1963}
E.~L.~Lawler,
``The quadratic assignment problem,''
\textit{Management Science},
vol.~9, no.~4, pp.~586--599, 1963.

\bibitem{leordeanu2005}
M.~Leordeanu and M.~Hebert,
``A spectral technique for correspondence problems using pairwise constraints,''
in \textit{Proc. IEEE Int. Conf. Computer Vision (ICCV)},
vol.~2, pp.~1482--1489, 2005.

\bibitem{birkhoff1946}
G.~Birkhoff,
``Tres observaciones sobre el algebra lineal,''
\textit{Universidad Nacional de Tucuman Revista Serie A},
vol.~5, pp.~147--151, 1946.

\bibitem{frank1956algorithm}
M.~Frank and P.~Wolfe,
``An algorithm for quadratic programming,''
\textit{Naval Research Logistics Quarterly},
vol.~3, no.~1--2, pp.~95--110, 1956.

\bibitem{jaggi2013revisiting}
M.~Jaggi,
``Revisiting Frank-Wolfe: Projection-free sparse convex optimization,''
in \textit{Proc. Int. Conf. Machine Learning (ICML)},
pp.~427--435, 2013.

\bibitem{kuhn1955hungarian}
H.~W.~Kuhn,
``The Hungarian method for the assignment problem,''
\textit{Naval Research Logistics Quarterly},
vol.~2, no.~1--2, pp.~83--97, 1955.

\bibitem{kabsch1976solution}
W.~Kabsch,
``A solution for the best rotation to relate two sets of vectors,''
\textit{Acta Crystallographica Section A},
vol.~32, no.~5, pp.~922--923, 1976.

\bibitem{schoenberg1938metric}
I.~J.~Schoenberg,
``Metric spaces and positive definite functions,''
\textit{Transactions of the American Mathematical Society},
vol.~44, no.~3, pp.~522--536, 1938.

\bibitem{shervashidze2011}
N.~Shervashidze, P.~Schweitzer, E.~J.~van~Leeuwen, K.~Mehlhorn, and K.~M.~Borgwardt,
``Weisfeiler-Lehman graph kernels,''
\textit{Journal of Machine Learning Research},
vol.~12, pp.~2539--2561, 2011.

\bibitem{aromatase2020}
A.~Ghisleni et al.,
``Targeting aromatase in breast cancer: Molecular mechanisms, therapeutic inhibitors, and clinical perspectives,''
\textit{Bioorganic \& Medicinal Chemistry},
vol.~28, no.~14, p.~115560, 2020.

\bibitem{peyre2019computational}
G.~Peyr{\'e} and M.~Cuturi,
``Computational optimal transport: With applications to data science,''
\textit{Foundations and Trends in Machine Learning},
vol.~11, no.~5--6, pp.~355--607, 2019.

\bibitem{titouan2019optimal}
T.~Vayer, L.~Chapel, R.~Flamary, R.~Tavenard, and N.~Courty,
``Optimal transport for structured data with application on graphs,''
in \textit{Proc. Int. Conf. Machine Learning (ICML)},
pp.~6275--6284, 2019.

\bibitem{vogelstein2015fast}
J.~T.~Vogelstein, J.~M.~Conroy, V.~Lyzinski, L.~J.~Podrazik, S.~G.~Kratzer, E.~T.~Flake, M.~W.~Fishkind, D.~E.~Fishkind, and C.~E.~Priebe,
``Fast approximate quadratic programming for graph matching,''
\textit{PLOS ONE},
vol.~10, no.~4, p.~e0121002, 2015.

\end{thebibliography}

\end{document}
